% ============================================================================
%  Subdocumento 4. Introducción a la Arquitectura Lógica y Física de la
%  Solución. Índice obligatorio del Comunicado 10, sección 11.
%
%  4.1 y 4.1.1 son de la dupla 3. Su texto se trasladó sin cambios del
%  Subdocumento 4 entregado en el Informe 1 (D4-51), salvo el párrafo sobre la
%  emisión a bordo del documento de transporte, que la observación 76 objetó.
%  4.2 a 4.3.2 son de la dupla 4 (decisiones D4-01 a D4-60).
% ============================================================================

Este capítulo describe cómo se construye la \dato{proyecto}: con qué piezas de software se organiza
la solución y en qué equipos, sitios y regiones de nube corre cada una. La sección~4.1 presenta la
arquitectura lógica, sus ocho capas, sus seis contextos delimitados y sus decisiones registradas, y la
4.1.1 las tecnologías de software. La sección~4.2 ubica cada componente en la nube, en San Bernardo,
en los cuatro terminales regionales y en los camiones, declara los enlaces y los puntos únicos de
falla, y dimensiona la ingesta de datos y la reconexión masiva. La sección~4.3 declara los dos
centros de datos: Azure Chile Central como región primaria, con la sala de San Bernardo como sitio de
continuidad, y Azure Brazil South como región secundaria.

El capítulo parte del esquema de solución del Subdocumento~3 y del catálogo de requerimientos de su
Formulario T-12. Se apoya en el modelo de datos del Subdocumento~5, que define dónde reside cada dato
y cómo se respalda. Los riesgos que abre la infraestructura, como el apagado de las redes 2G y 3G o el
retraso de una réplica, se tratan en el Subdocumento~8, y el ritmo de instalación a bordo en el plan
de trabajo del Subdocumento~7. La innovación tecnológica del Subdocumento~13 usa el Bluetooth del
equipo a bordo que se especifica aquí. Lo que el mandante compra, con cantidades y características, se
detalla en el Formulario T-11, que acompaña este capítulo como archivo propio
(\texttt{AUDIT-Formulario-T-11.pdf}), conforme al \comunicado{10}{sección 1}.

% ============================================================================
\section{Arquitectura lógica}\label{sec:4-logica}\label{sec:04-arquitectura-logica}

El Capítulo 5 del Caso resume el problema en una frase. El sistema de gestión de transporte de 2013 conoce el viaje que la compañía encargó y no conoce el viaje que efectivamente ocurrió (\caso{capítulo 5}{13}). El dato existe repartido en tres plataformas de posicionamiento, en una telemetría que nadie descarga, en una liquidación de combustible que llega con cuarenta días de atraso y en papeles que viajan en la cabina, y nunca se junta. Esa dispersión no se corrige configurando el sistema existente, porque nace de cómo está construido.

\subsection{Especificaciones Tecnologías de Software a utilizar}\label{sec:4-software}\label{sec:04-especificaciones-tecnologias-de-software-a-utilizar}

La selección sigue RT-03.05: servicios administrados cuando reducen el riesgo operacional, con separación entre datos de negocio, series de tiempo y análisis. La \tabref{04-tabla-4-1} fija productos y políticas de soporte. El contrato de 56 meses exige actualizaciones; una versión inicial no se presenta como soportada hasta el mes 56.

\begin{tabla}[ancho=completo, fuente=condensada]{Tecnologías seleccionadas y mantenimiento del soporte}{04-tabla-4-1}{L{36mm} X X}
Componente & Selección y fundamento & Soporte y actualización \\ \midrule
Aplicación móvil, cuatro perfiles & Flutter 3, Dart y SQLite local: una base de interfaz y operación desconectada. Se descartan interfaces exclusivamente web para cabina por integración local y pruebas con guantes & Canal estable; fijar versión exacta y plataformas mínimas en el manifiesto del hito de diseño. No se atribuye un LTS indefinido a Flutter \\
Portales y torre & React y Next.js; separar consulta pública de estado y portales autenticados. WCAG 2.2 AA verificable en S9 & Versiones exactas, licencias y dependencias en SBOM; actualización y regresión antes del fin de soporte de la versión escogida \\
Servicios de negocio & .NET 10 LTS; Go 1.27 para ingesta/borde. Se conserva separación de contextos y contratos & .NET 10 hasta 14-11-2028; Go mantiene cada versión hasta dos versiones mayores posteriores. .NET 8 termina 10-11-2026 y no es la base del inicio supuesto en febrero de 2027 \\
Orquestación & Azure Kubernetes Service; despliegue independiente y aislamiento de recursos & Versión soportada por AKS en cada liberación; matriz exacta de Kubernetes e imágenes, actualización por grupos y pruebas de reversa \\
Transaccional & PostgreSQL 16 Flexible Server con PostGIS, conforme a T-11; claves y restricciones para reservas y vigencias & PostgreSQL 16 hasta 09-11-2028. Ensayar actualización mayor antes de esa fecha; preservar compatibilidad de extensiones y CDC \\
Series de tiempo & Instancia Flexible Server separada con TimescaleDB, edición Apache-2 disponible en Azure. Se descarta Cosmos DB para este dominio por conservar SQL y el modelo temporal elegido & Azure publica la matriz de versiones de extensión por versión de PostgreSQL. No se inventa un fin de soporte independiente en mayo de 2029 ni se presuponen funciones comerciales de compresión no incluidas \\
Eventos & Azure Event Hubs Premium, compatible con clientes del protocolo Kafka; se descarta operar un clúster Strimzi adicional para el mismo flujo & Servicio administrado. Orden por partición y entrega al menos una vez; idempotencia propia. No depender de transacciones Kafka en vista previa \\
Caché & Azure Managed Redis, producto de T-11; almacena datos derivados y sesiones, nunca la única prueba habilitante de despacho & Para una implantación nueva se provisiona el destino. Sólo si hay una instancia de Azure Cache for Redis efectivamente existente se programa su migración antes del retiro de su SKU \\
Persistencia local & SQLite 3 en WAL; cola de salida transaccional y acuses duraderos & Fijar biblioteca e imagen Linux del equipo; probar cortes de energía y actualizar sin borrar registros no confirmados. No se denomina LTS al producto sin política aplicable \\
Analítica & ADLS Gen2 y tablas Delta Lake; cómputo y modelo dimensional separados del OLTP & Manifiesto de versiones del motor de procesamiento, conectores y biblioteca Delta; exportación en formatos abiertos \\
Identidad y secretos & Entra ID, identidades administradas y Key Vault; autorización por rol, atributo, transportista y viaje & Fijar edición y capacidades en T-11; no atribuir certificación de HSM a una bóveda o SKU sin comprobarlo \\
\end{tabla}

La \tabref{04-tabla-4-1} decide productos y evita dos buses alternativos o dos motores intercambiables. Las versiones exactas de los canales móviles y web se fijan en el hito de diseño con SBOM, licencia, prueba de compatibilidad, fecha o política de soporte y responsable. Esta es una especificación de mantenimiento, no una prueba de concepto ya ejecutada. Los servicios cloud y su disponibilidad regional se cotejan con S4.2/T-11 antes de su aprovisionamiento.

Fuentes consultadas el 11-10-2026: política de .NET \parencite{d3_net}, política de Go \parencite{d3_go}, soporte de PostgreSQL \parencite{d3_postgres}, extensiones de Azure \parencite{d3_extensiones}, consideraciones de TimescaleDB \parencite{d3_timescale} y protocolo Kafka en Event Hubs \parencite{d3_kafka}.

Microsoft lista tres zonas en Chile Central y no le asigna región emparejada (regiones de Azure \parencite{d3_regiones}). Brazil South es una elección de DR de la propuesta. La existencia de la región no acredita disponibilidad de cada SKU, extensión, modalidad HA o cuota: D3/D4 deben fijar la matriz por servicio y región en el diseño detallado y validar aprovisionamiento y recuperación antes de activar producción.

\subsection{Punto de partida}\label{sec:4-partida}\label{sec:04-punto-de-partida}

Los rasgos del sistema de 2013 que explican esa dispersión y la respuesta que da esta arquitectura a cada uno se resumen en la \tabref{04-tabla-4-2}.

\begin{tabla}[ancho=completo, fuente=condensada]{Rasgos del sistema de 2013 y respuesta de esta arquitectura}{04-tabla-4-2}{L{36mm} X X}
Rasgo & Efecto que produce hoy & Respuesta de esta arquitectura \\ \midrule
Base operativa del sistema de gestión 2013 acoplada a consultas de gestión & La consulta de gestión compite con la operación de la torre por el mismo motor & Separación del almacenamiento transaccional y del analítico, obligatoria por RT-05.05 \\
Integración con el sistema contable por consultas y tablas compartidas & Todo cambio contable alcanza el núcleo operacional sin traducción & Capa anticorrupción, obligatoria por RT-05.20 \\
Ausencia de costo por kilómetro por ruta & Tres de los ocho contratos principales se sirven bajo costo, el peor a menos catorce por ciento sostenido durante cuatro años (\caso{numeral 7.3}{15}) & Costo por viaje en doble versión, preliminar y consolidada \\
Puerto de telemetría de fábrica inactivo & La lectura del motor de 61 tractocamiones no se aprovecha por temor a la garantía & Lectura de fábrica autorizada; protocolo y equipo homologados con cada fabricante \\
Módulos operativos acoplados entre sí & El sistema no distingue el viaje encargado del viaje ocurrido & Seis contextos delimitados, cada uno con sus propios datos \\
\end{tabla}

\notaTabla{Fuente: Bases Técnicas del Caso, capítulos 5 y 7.}

Cada rasgo de la tabla tiene una respuesta estructural propia. Ninguna se resuelve configurando el sistema existente.

\subsection{Las ocho capas}\label{sec:4-capas}\label{sec:04-las-ocho-capas}

La solución se organiza en las ocho capas del modelo de referencia, cuya existencia es obligatoria conforme al \transv{numeral 2.1}{6}. RT-02.01 exige el diagrama que identifica cada capa, sus componentes y las interfaces entre ellas, y la descripción se ajusta a ISO/IEC/IEEE 42010 \parencite{iso42010}. La \tabref{04-tabla-4-3} resume el contenido de cada capa.

\begin{tabla}[ancho=completo, fuente=condensada]{Las ocho capas y su contenido en esta solución}{04-tabla-4-3}{L{47mm} X}
Capa & Contenido \\ \midrule
Presentación & Portal web para 84 clientes y 148 transportistas, aplicación móvil, terminales de torre y taller, pantallas de terreno operables con guantes \\
Borde y exposición & Distribución de contenidos, cortafuegos de aplicación, balanceo y terminación de cifrado. Único punto de entrada público \\
Puerta de enlace & Autenticación, autorización, cuotas, límites de tasa, versionado semántico y catálogo de servicios \\
Servicios de negocio & Seis contextos delimitados con límites explícitos, sin estado y desplegables de forma independiente \\
Integración y eventos & Bus de telemetría, bus transaccional, cola de mensajes fallidos, reintento y deduplicación \\
Datos & Persistencia políglota. Transaccional, series de tiempo, documental inmutable y analítica separadas \\
Seguridad & Identidad, autorización, gestión de secretos, cifrado y auditoría. Transversal a todas las capas \\
Observabilidad & Métricas, registros y trazas correlacionadas, sin puntos ciegos entre nube y terreno \\
\end{tabla}

\notaTabla{Fuente: Bases Técnicas Transversales, numeral 2.1, y elaboración propia.}

Ninguna interfaz accede directamente a la base de datos. Una petición entra por presentación, borde y puerta de enlace; el servicio invoca su repositorio o publica mediante outbox cuando corresponde. No toda consulta atraviesa el bus de eventos. La respuesta vuelve por la puerta de enlace y el canal de presentación. Seguridad y observabilidad no ocupan un lugar en esa pila, la atraviesan entera. Los componentes declarados en cada capa se listan en la \tabref{04-tabla-4-4}.

\begin{tabla}[ancho=completo, fuente=condensada]{Componentes declarados por capa}{04-tabla-4-4}{L{47mm} X}
Capa & Componentes \\ \midrule
Presentación & Portal web con representación en servidor, aplicación móvil en los cuatro perfiles que exige RT-17.01 del Caso, conductor, torre, terminal y taller, y transportista subcontratado, y vistas de portería y de taller de alto contraste \\
Borde y exposición & Distribución de contenidos con presencia global, cortafuegos de aplicación y protección volumétrica en las capas de red, transporte y aplicación \\
Puerta de enlace & Gestión de interfaces con identidad federada, certificado mutuo entre máquinas, límites de tasa por perfil y catálogo publicado \\
Servicios de negocio & Contenedores orquestados sobre nodos repartidos en tres zonas de disponibilidad, con un despliegue independiente por contexto delimitado \\
Integración y eventos & Flujo de telemetría para la ingesta masiva, bus transaccional con orden garantizado dentro de la partición y cola de mensajes fallidos, y capa anticorrupción frente al sistema contable \\
Datos & Motor relacional multizona, base de series de tiempo, caché en memoria, almacenamiento inmutable y repositorio analítico. El diseño detallado es materia del Subdocumento 5 \\
Seguridad & Bóveda de claves en módulo criptográfico, identidad federada con control por rol y por atributo, y bitácora que permite reconstruir quién, qué, cuándo y con qué valores anteriores, según RT-05.03 \\
Observabilidad & Instrumentación única con trazas, métricas y registros correlacionados por el identificador común que RT-05.19 exige a toda integración \\
\end{tabla}

\notaTabla{Fuente: elaboración propia.}

La \figref{04-logicacapas} dibuja las ocho capas con sus componentes.

\anexoGrafico{../../Informe-2/D3/graficos/LogicaCapas.png}{Las ocho capas obligatorias del numeral 2.1 transversal, con sus componentes e interfaces, conforme a RT-02.01}{04-logicacapas}

Fuente: elaboración de audIT.

La figura muestra la pila de seis capas y las dos capas transversales. Las capas de servicios y de datos existen además a bordo del camión, en versión reducida, para operar sin enlace.

\subsection{Los seis contextos delimitados}\label{sec:4-contextos}\label{sec:04-los-seis-contextos-delimitados}

RT-02.02 rechaza toda arquitectura monolítica que no permita desplegar de forma independiente sus componentes críticos. Los servicios de negocio se organizan en seis contextos, cada uno con su propio lenguaje y sus propios datos. Un contexto no consulta la base de datos de otro, le pregunta. La \tabref{04-tabla-4-5} declara qué decide cada uno.

\begin{tabla}[ancho=completo, fuente=condensada]{Contextos delimitados y la decisión que sostiene cada uno}{04-tabla-4-5}{L{47mm} X}
Contexto & Qué decide \\ \midrule
Planificación y tráfico & Qué carga hay que mover y hacia dónde \\
Flota y activos & Si el equipo puede salir hoy \\
Personas y cumplimiento & Si la persona puede conducir hoy \\
Telemetría y geocercas & Dónde está el camión y cuándo llegó \\
Operación de fletes & El viaje que realmente ocurrió \\
Liquidación y costeo & Cuánto costó y a quién se paga \\
\end{tabla}

\notaTabla{Fuente: elaboración propia.}

La descomposición separa responsabilidades que hoy quedan dispersas entre el sistema de gestión, planillas, papeles y plataformas de posicionamiento. La \figref{04-audit-figura-4-2-contextos} muestra los seis contextos y los sistemas con los que conviven.

\anexoGrafico{../../Informe-2/D3/graficos/audit-figura-4.2-contextos.png}{Contextos delimitados del dominio y sistemas con los que convive la solución}{04-audit-figura-4-2-contextos}

Fuente: D4, figura conservada.

En la figura, cada contexto tiene su propio lenguaje y sus propios datos. Planificación y tráfico puede bloquear un viaje sin conocer por dentro a Personas y cumplimiento ni a Flota y activos, porque les pregunta.

\subsection{Convivencia con lo que ya existe}\label{sec:4-convivencia}\label{sec:04-convivencia-con-lo-que-ya-existe}

El sistema de gestión de transporte de 2013 y el sistema contable son contrapartes distintas (\caso{capítulo 5}{12--13}). La capa anticorrupción traduce sus contratos, sin sustituir por sí sola sus módulos. Los servicios nuevos reemplazan asignación, control de viaje y liquidación en M16; órdenes y tarifas a clientes en M21. El sistema de gestión queda de consulta hasta M24, con conciliación y reversa según S3/S7. El contable permanece como único emisor tributario bajo la Restricción 8.

\subsection{Registro de decisiones de arquitectura}\label{sec:4-adr}\label{sec:04-registro-de-decisiones-de-arquitectura}

RT-02.04 exige ADR fechados y actualizados. La \tabref{04-tabla-4-6} registra ocho decisiones de diseño al 11-10-2026. Se preservan los identificadores de las cuatro decisiones previas para no romper las referencias de S4.2; los nuevos ADR 05–08 completan el análisis solicitado para el Informe 2.

\begin{tabla}[ancho=completo, fuente=condensada]{Registro de decisiones de arquitectura}{04-tabla-4-6}{L{29mm} X X X}
ADR / fecha & Decisión & Alternativas evaluadas & Criterio, consecuencia y verificación \\ \midrule
01 / 11-10-2026 & Seis contextos desplegables en AKS & Monolito modular, SOA con orquestador central, servicios independientes & RT-02.02 exige independencia de componentes críticos. Separar ingesta del despacho; asumir mayor costo de operación, observabilidad y pruebas de contratos. Medir aislamiento y carga en S9 \\
02 / 11-10-2026 & Event Hubs Premium con protocolo Kafka y outbox transaccional & Cadena HTTP síncrona; Kafka autoadministrado sobre Strimzi & Aislar ráfagas y evitar que los efectos derivados alarguen el despacho. Deduplicar y reconciliar entrega al menos una vez; no confundir el endpoint Kafka con un clúster Kafka propio \\
03 / 11-10-2026 & CDC hacia ADLS/Delta y modelo dimensional & BI sobre OLTP; réplica operacional usada para BI & RT-05.05 exige separación. CDC también consume recursos: vigilar slots/WAL y medir su impacto; no prometer carga nula \\
04 / 11-10-2026 & Identidad federada y certificados por equipo; acceso por rol y atributos & Claves compartidas o parámetros secretos en URI & RT-05.18 y Zero Trust: revocación, mínima exposición, secreto fuera del código y autorización por viaje. Probar acceso cruzado entre transportistas \\
05 / 11-10-2026 & Reemplazar gestión 2013 función por función; conservar contable & Corte único; conservar sólo integrando; paquete de mercado sin integración demostrada & M16/M21/M24 de S3. Evitar interrupción y doble registro definitivo; presupuestar coexistencia, adaptadores y conciliación en los paquetes de S7 y Sobre 3, sin inventar ahorro monetario \\
06 / 11-10-2026 & Captura local y enlace satelital de mensajes; documento emitido en contable & Banda ancha satelital en toda la flota; generación fiscal en cabina & Restringir costo recurrente y equipos del T-11; la alternativa local infringe R-08. El descarte de banda ancha no prueba que SBD baste: validar entrega documental completa y 90 s; revisar esta decisión si el ensayo falla \\
07 / 11-10-2026 & PostgreSQL transaccional y TimescaleDB separado & Motor único; Cosmos DB; motor temporal autoadministrado adicional & Integridad referencial en el negocio, aislamiento de series y SQL conocido. Verificar extensión Apache-2, funciones efectivas, actualización y HA en Azure \\
08 / 11-10-2026 & SQLite local con outbox, acuses y conciliación determinista & Depender del enlace para cada evento; búfer sólo en RAM & 72 h a bordo, 24 h en terminal como mínimo transversal y confirmación $\leq$20 min tras 72 h. 288 h es ampliación propuesta: ensayos y capacidad separados \\
\end{tabla}

La \tabref{04-tabla-4-6} permite revisar la elección y sus costos operativos sin declarar pruebas ejecutadas. Cada cambio de proveedor, protocolo o estrategia de continuidad actualiza el ADR, los contratos, la EDT y los casos de aceptación.

\subsection{Degradación, escalamiento y puntos únicos de falla}\label{sec:4-degradacion}\label{sec:04-degradacion-escalamiento-y-puntos-unicos-de-falla}

Tres requisitos obligatorios de la misma sección exigen declaraciones que conviene no dejar implícitas. RT-02.09 obliga a degradar de forma elegante, de modo que la caída de un componente no crítico deje la solución operando en modo reducido y avisando de la degradación, y nunca falle de forma total. La aplicación de esa regla en esta operación es directa. Si el repositorio analítico no responde, la torre sigue despachando. Si el sistema contable no responde, continúan la captura, el seguimiento y los viajes que ya disponen de documentación conforme. Se encola la solicitud pendiente, pero se bloquea la salida de un viaje que todavía requiere emisión. Si el enlace satelital no está disponible, el registro sigue acumulándose a bordo.

RT-02.10 exige que las capas de aplicación e integración escalen horizontalmente de forma automática, con umbrales, límites superiores y costo asociado declarados en la oferta. Los umbrales y los límites se declaran en esta sección. El costo asociado se remite al Sobre N.º 3, porque el \art{50.2}{29} excluye toda cifra de precio de la Oferta Técnica.

RT-02.11 evalúa como observación grave omitir la declaración de los puntos únicos de falla que subsistan. Esta oferta declara dos en la capa lógica, que se listan en la \tabref{04-tabla-4-7}.

\begin{tabla}[ancho=completo, fuente=condensada]{Puntos únicos de falla declarados en la capa lógica}{04-tabla-4-7}{L{36mm} X X}
Punto & Por qué subsiste & Por qué es aceptable \\ \midrule
El sistema contable como emisor único de documentos tributarios & La restricción 8 lo fija y no es negociable & La captura y el seguimiento continúan. Una salida sin documento conforme queda bloqueada; se monitoriza la cola y se escala como incidente crítico \\
El dispositivo a bordo de cada camión & Hay uno por unidad y solo puede intervenirse cuando el camión pasa por un terminal, según RT-06.01 del Caso & La falla afecta a una unidad y no a la flota. Se mitiga con repuestos precargados y con el procedimiento de contingencia aprobado, sin levantar bloqueos legales \\
\end{tabla}

\notaTabla{Fuente: elaboración propia sobre el Caso, capítulo 10 y RT-06.01.}

Los puntos únicos de falla de la infraestructura física se analizan en la sección 4.2.7.

\subsection{Resiliencia}\label{sec:4-resiliencia}\label{sec:04-resiliencia}

Los servicios de negocio son sin estado, con el estado de sesión y de proceso en almacenes externos de alta disponibilidad. Los patrones de RT-02.08 son obligatorios y esta oferta los adopta. Tiempos límite explícitos en toda llamada remota, cortacircuitos y mamparos para aislar fallas de integraciones externas, y reintento exponencial con variación aleatoria. La escritura es idempotente con ventana de deduplicación dimensionada para tolerar la desconexión prolongada en ruta.

El estado de sesión y el estado de proceso residen en almacenes externos de alta disponibilidad, como exige RT-02.05. Allí viven las sesiones concurrentes de los 22 operadores de la torre, las credenciales activas y las geocercas de los 1.400 puntos distintos de carga y descarga que declara el \caso{numeral 14.1}{29}. Cualquier instancia puede destruirse y reemplazarse sin pérdida de transacciones en curso. Los parámetros de la \tabref{04-tabla-4-8} son de diseño de audIT y se ajustan con la medición de la Etapa 1.

\begin{tabla}[ancho=completo, fuente=condensada]{Parámetros de los patrones de resiliencia}{04-tabla-4-8}{L{47mm} X}
Patrón & Parámetro declarado \\ \midrule
Cortacircuito & Se abre al superarse la mitad de las llamadas fallidas en una ventana deslizante de diez peticiones, permanece abierto 30 segundos y admite tres llamadas de prueba antes de restablecer el tráfico \\
Mamparo & La verificación bloqueante del despacho tiene su propio conjunto de hilos y de conexiones. La saturación de la consulta de liquidaciones o de reportes no le quita capacidad \\
Tiempo de espera & Validación en memoria, 800 milisegundos. Consulta transaccional del despacho, 5 segundos. Llamada al sistema contable a través de la capa anticorrupción, 10 segundos. Emisión del documento electrónico de transporte, 90 segundos, que es el techo que fija RT-09.01 del Caso \\
Reintento & Hasta tres intentos sobre operaciones idempotentes sólo si caben en el plazo total, con retroceso exponencial y variación aleatoria. La variación evita que cientos de unidades que recuperan cobertura a la vez reintenten sincronizadas \\
Límite de tasa & Declarado por perfil en la puerta de enlace, según se detalla en la gobernanza de interfaces \\
\end{tabla}

\notaTabla{Fuente: elaboración propia sobre RT-02.08 y RT-09.01 del Caso.}

Los tiempos de espera escalonan desde la validación en memoria hasta el techo de 90 segundos del Caso, de modo que ninguna llamada lenta consume el presupuesto de la verificación bloqueante.

\subsection{Verificación bloqueante del despacho}\label{sec:4-verificacion}\label{sec:04-verificacion-bloqueante-del-despacho}

\caso{RT-09.01}{32} fija en 30 segundos el tope para asignar un viaje con verificación de jornada, habilitaciones y aptitud del equipo. El servicio evalúa en paralelo las tres invariantes de la \tabref{04-tabla-4-9}, y ninguna admite excepción automática.

\begin{tabla}[ancho=completo, fuente=condensada]{Invariantes de la verificación bloqueante}{04-tabla-4-9}{L{36mm} X X}
Invariante & Comprobación & Fuente y efecto \\ \midrule
Persona y jornada & Identidad del conductor, descanso y jornada acumulada, evidencia previa relevante y licencia vigente & S3 RN-01–04; \caso{RT-09.01}{32}. Una causa legal bloquea siempre \\
Equipo & Vigencias acreditadas de tracto y semirremolque, disponibilidad y compatibilidad con la carga & FEP03 · numerales 4.4 y 16.1; S3 RN-05. Sin respaldo vigente no habilita \\
Carga y documentación & Documento de transporte conforme, requisitos y lista de verificación firmada para carga peligrosa & Restricción 8; RF-006/RF-013/RF-014. La asignación no sustituye la autorización de salida \\
\end{tabla}

La \tabref{04-tabla-4-9} separa la elegibilidad de los recursos y la condición documental de salida. El conductor, tracto y semirremolque no pueden tener asignaciones activas superpuestas. El servicio obtiene evaluaciones versionadas de Personas y de Flota; Planificación confirma únicamente si esas versiones siguen vigentes. La autorización de salida exige además el documento y la comprobación de carga.

Cada evaluación lleva sujeto, versión, instante y vencimiento. No basta leer una versión remota y después hacer commit en otra base: la revalidación y el protocolo de reserva deben impedir confirmar mientras cambia una habilitación. Planificación serializa las reservas de recursos y conserva el intento antes de consultar. Una modificación incompatible invalida los intentos afectados; ante falta de confirmación de vigencia se rechaza dentro del presupuesto total. S9 prueba especialmente una revocación concurrente, dos solicitudes para el mismo recurso y el vencimiento del token; no se afirma atomicidad global por usar una transacción local.

La \tabref{04-tabla-4-10} distribuye el tiempo máximo de la decisión, incluyendo confirmación, auditoría y respuesta.

\begin{tabla}[ancho=completo, fuente=condensada]{Presupuesto de diseño de la asignación, extremo a extremo}{04-tabla-4-10}{L{36mm} X X}
Paso & Máximo presupuestado & Control \\ \midrule
Identidad, contrato y clave de idempotencia & 2 s & Autorización y unicidad duradera \\
Jornada, vigencias y aptitud, en paralelo & 8 s & Evidencia vigente; retorno a fuente autoritativa si caché desactualizada \\
Reserva y confirmación atómica & 8 s & Restricciones de exclusión/unicidad, versión y conflicto de concurrencia \\
Composición del veredicto y auditoría & 4 s & Causal, evidencia, regla y correlación \\
Respuesta al operador & 3 s & Incluye tránsito de red y puerta de enlace \\
Total + margen & 25 s + 5 s = 30 s & Ningún reintento rebasa el plazo total \\
\end{tabla}

La \tabref{04-tabla-4-10} es una asignación de tiempo de diseño, no una medición. Los 660 ms nominales representados en la \figref{04-audit-figura-4-3-resiliencia} describen un caso favorable de caché; no acreditan extremo a extremo ni reemplazan este presupuesto. Un objetivo interno menor de 2 s se conserva como hipótesis de optimización y se mide en S9, sin presentarlo como resultado demostrado.

Ante rechazo, la respuesta identifica causa, regla aplicada, valor evaluado, dato faltante y correlación; la \tabref{04-tabla-4-11} fija los campos. Una petición repetida devuelve el resultado ya confirmado y una operación aún en curso devuelve conflicto, sin emitir una segunda asignación.

\begin{tabla}[ancho=completo, fuente=condensada]{Respuesta trazable ante rechazo}{04-tabla-4-11}{L{47mm} X}
Campo & Contenido \\ \midrule
tipo / estado & Causal estable y categoría: regla, dato insuficiente, conflicto o error técnico \\
detalle & Umbral, dato evaluado y motivo, sin divulgar antecedentes personales a un rol no autorizado \\
instancia / correlacion & Viaje, intento y operaciones relacionadas \\
versionRegla / evidencias & Versión normativa-operacional y referencias que sustentan el veredicto \\
\end{tabla}

La \tabref{04-tabla-4-11} permite al operador corregir la causa y conserva la decisión para auditoría. RN-04 permite una excepción sólo por fuente indisponible, con jefe de turno y prevención de riesgos, motivo, un viaje y revisión mensual; no admite jornada agotada ni documento legal faltante.

La \figref{04-audit-figura-4-3-resiliencia} muestra los patrones que aíslan el despacho de las integraciones externas.

\anexoGrafico{../../Informe-2/D3/graficos/audit-figura-4.3-resiliencia.png}{Patrones de resiliencia y presupuesto nominal de diseño}{04-audit-figura-4-3-resiliencia}

La \figref{04-4-3-asignacion} detalla los mensajes de la asignación. La transacción incorpora reserva, asignación y outbox; sólo después del commit se devuelve confirmación. La publicación se reintenta desde outbox y el consumidor deduplica. Se elimina así la ventana de pérdida entre confirmar el viaje y publicar su evento.

\anexoGrafico{../../Informe-2/D3/graficos/4-3-asignacion.png}{Secuencia de asignación bloqueante y publicación con outbox}{04-4-3-asignacion}

La rama de rechazo de la \figref{04-4-3-asignacion} conserva el bloqueo ante incumplimiento o plazo agotado. La excepción RN-04 es un procedimiento separado y aprobado, nunca un reintento que cambia silenciosamente el veredicto.

\subsection{Idempotencia y ventana de deduplicación}\label{sec:4-idempotencia}\label{sec:04-idempotencia-y-ventana-de-deduplicacion}

RT-02.06 exige escrituras idempotentes. Cada operación lleva una clave única. La caché puede retenerla siete días para acelerar reintentos frecuentes; siete días no cubren doce días. La unicidad duradera y la respuesta registrada permanecen durante la retención del registro de negocio, incluida cualquier reconexión de 288 h. La clave se bloquea en el almacén en memoria y se respalda con una restricción de unicidad duradera en el motor transaccional, de modo que un reintento tras la reconexión masiva no duplica un viaje ni un documento.

La ventana opera en dos niveles con propósitos distintos. El nivel en memoria resuelve la concurrencia inmediata. Si la clave ya está tomada y la operación sigue en curso, la petición se rechaza como conflicto. Si ya se completó, se devuelve la respuesta guardada en lugar de repetir la escritura. El nivel duradero resuelve el caso que importa aquí, el mensaje que llega desde el búfer de un camión después de días sin cobertura, cuando la clave ya expiró en memoria. La restricción de unicidad lo descarta sin importar cuánto tiempo pasó. RT-02.07 exige además deduplicación en el consumidor y orden garantizado dentro de la partición, y ambos se aplican al flujo de telemetría.

\subsection{Gobernanza de interfaces y contratos}\label{sec:4-interfaces}\label{sec:04-gobernanza-de-interfaces-y-contratos}

\transv{RT-05.16}{12} obliga a documentar los servicios síncronos en OpenAPI 3.1 y los flujos dirigidos por eventos en AsyncAPI 2.6 o superior, y exige que esa documentación se genere desde el código y se mantenga actualizada de forma automática. La consecuencia práctica es que no se admite discrepancia entre el contrato publicado y la implementación viva, porque el contrato no se escribe aparte. A modo representativo, el servicio síncrono de verificación bloqueante pre-despacho expone el contrato \texttt{POST /api/v1/despacho/validar} documentado bajo OpenAPI 3.1, recibiendo en su esquema JSON los identificadores del viaje (\texttt{viajeId}), conductor (\texttt{conductorId}), tracto (\texttt{patenteTracto}), rampla (\texttt{patenteRampla}) y carga (\texttt{codigoSUSPEL}), retornando el veredicto (\texttt{autorizado}), latencia de validación (\texttt{latenciaMs}) y la bitácora auditable de factores. En paralelo, el evento asíncrono de notificación de salida se formaliza bajo AsyncAPI 2.6 en el canal de eventos de operaciones, transportando eventos CloudEvents con identificador único, agregado, secuencia, versión de esquema y correlación. El transporte autenticado y el sellado de evidencias se especifican según el tipo de evento; un HMAC no se presenta como firma tributaria.

RT-05.17 obliga a versionar semánticamente los contratos, con compatibilidad hacia atrás y política de obsolescencia con preaviso mínimo de seis meses. RT-05.18 fija el mecanismo de autenticación entre sistemas y prohíbe la clave estática en la ruta de la dirección web. RT-05.19 exige registrar la transacción de entrada y la de salida de toda integración con un identificador de correlación común, que es lo que permite seguir una operación de negocio a través de todos los sistemas que toca. Ese identificador es el mismo que usa la capa de observabilidad.

RT-05.21 obliga a declarar, por cada integración, el modo, el volumen esperado, la ventana de disponibilidad de la contraparte y el comportamiento de la solución cuando esa contraparte no responde. Esa declaración se entrega en la \tabref{04-tabla-4-12}.

\begin{tabla}[ancho=completo, fuente=condensada]{Declaración por integración conforme a RT-05.21}{04-tabla-4-12}{L{27mm} X X X X}
Contraparte & Modo & Volumen esperado & Ventana de la contraparte & Si no responde \\ \midrule
Sistema contable y de facturación & Solicitud/respuesta para emisión; asíncrono para asientos & Hasta 128.000 documentos y eventos de 96.000 viajes/año & Capacidad 24x7 de emisión por acordar y probar; no consta en el Caso & Cortacircuito y encolado. Sin documento conforme se bloquea la salida; captura y viajes ya documentados continúan \\
Plataformas de posicionamiento de terceros & Asíncrono & Posición de los camiones de terceros con dispositivo & No declarada por el Caso & Última posición conocida con su antigüedad visible, nunca una posición sin fecha \\
Telemetría de fábrica & Asíncrono & 61 tractocamiones, solo lectura & Sujeta a la autorización de cada fabricante & El viaje se registra igual con el dispositivo a bordo \\
Red de estaciones de servicio & Por lotes & 74.000 abastecimientos al año & Mensual, con hasta 40 días de desfase & El costo se emite preliminar y se marca como pendiente \\
Concesionarias de peaje & Por lotes & 620.000 pasadas al año & Mensual & Peaje estimado por la traza, marcado como estimado \\
Autoridad tributaria, a través del contable & Según contrato fiscal del emisor & 128.000 documentos/año & Según servicio de la autoridad & Aplicar contingencia fiscal admitida y documentada; audIT no se convierte en emisor \\
Clientes y autoridad aduanera & Ambos & Según contrato y cruce & Variable & Reintento con retroceso y aviso a la torre \\
\end{tabla}

\notaTabla{Fuente: \caso{numeral 14.1}{29}, y elaboración propia.}

Los volúmenes de esta tabla provienen del \caso{numeral 14.1}{29}. La ventana de disponibilidad de las plataformas de terceros y de la autoridad tributaria no está declarada en las bases y se consulta al mandante. RT-05.22 exige además carga y descarga masiva en formatos abiertos, con validación previa, informe de errores por registro y procesamiento parcial. Esa capacidad es la que sostiene la migración de las cerca de 6.000 vigencias y la ingesta mensual de combustible y peajes. La \figref{04-logicaintegraciones} dibuja el mapa de integraciones.

\anexoGrafico{../../Informe-2/D3/graficos/LogicaIntegraciones.png}{Mapa de integraciones. Sistemas internos, fuentes de terreno y contrapartes externas}{04-logicaintegraciones}

Fuente: D4, figura conservada.

La \figref{04-logicaintegraciones} diferencia el sistema de gestión 2013, que se sustituye, del contable, que se conserva. Los 192 terceros con equipo son población de referencia; el dibujo no acredita acceso automático a las tres plataformas. La cobertura efectiva y su antigüedad se muestran por unidad.

\subsection{Convivencia con el sistema contable heredado}\label{sec:4-contable}\label{sec:04-convivencia-con-el-sistema-contable-heredado}

Los servicios de negocio sustituyen los módulos operativos del sistema de gestión 2013. La capa anticorrupción traduce los contratos y encapsula el sistema contable, que se conserva como único emisor de documentos tributarios. El aislamiento es bidireccional. Un cambio en el sistema externo no propaga su modelo al núcleo, y el núcleo no escribe directamente sobre el legado.

RT-02.14 valora la aplicación documentada de patrones de arquitectura evolutiva que permitan sustituir un componente sin reescribir la solución, y nombra tres. Esta oferta aplica los tres. La capa anticorrupción frente al sistema heredado, el estrangulamiento progresivo con que los módulos operativos de 2013 se van sustituyendo función por función en lugar de en un corte único, y la abstracción de proveedores, que es lo que permite declarar la estrategia de reversibilidad que exige RT-03.07. Las piezas de la capa se listan en la \tabref{04-tabla-4-13}.

\begin{tabla}[ancho=completo, fuente=condensada]{Piezas de la capa anticorrupción}{04-tabla-4-13}{L{47mm} X}
Pieza & Función \\ \midrule
Adaptador de dominio & Traduce el evento de negocio, viaje cerrado o liquidación aprobada, a la estructura plana que el sistema contable espera \\
Transformador de esquemas & Homologa formatos en ambos sentidos, de modo que las convenciones de 2013 no lleguen al modelo nuevo \\
Protector de resiliencia & Cortacircuito y cola de mensajes fallidos. Si el contable no responde, las solicitudes quedan pendientes y las salidas sin documento conforme se bloquean \\
Reconciliador & Verifica que todo evento encolado terminó registrado, y expone el pendiente en lugar de dejarlo silencioso \\
\end{tabla}

\notaTabla{Fuente: elaboración propia.}

La \figref{04-audit-figura-4-4-acl} muestra la integración con el sistema contable a través de esa capa.

\anexoGrafico{../../Informe-2/D3/graficos/audit-figura-4.4-acl.png}{Integración con el sistema contable heredado a través de la capa anticorrupción}{04-audit-figura-4-4-acl}

Fuente: D4, figura conservada.

El sistema contable queda detrás de la capa y solo recibe eventos ya traducidos.

\subsection{Emisión del documento de transporte bajo Restricción 8}\label{sec:4-det-restriccion8}\label{sec:04-emision-del-documento-de-transporte-bajo-restriccion-8}

El sistema contable y de facturación es el único emisor fiscal (\caso{capítulo 10 · Restricción 8}{24}). Cabina y terminal no generan folios ni firman documentos tributarios. RF-013/RF-014 exigen documento conforme antes de iniciar el traslado; asignar recursos no habilita por sí solo la salida.

La \figref{04-documento-sin-cobertura} representa los dos caminos de S3: emisión anticipada cuando la carga definitiva se conoce, y solicitud satelital cuando sólo se conoce en el punto sin cobertura. En ambos, el contable emite y se entrega a cabina el contenido íntegro necesario para portar una representación conforme; un folio o una huella no bastan.

\anexoGrafico{../../Informe-2/D3/graficos/documento-sin-cobertura.png}{Emisión anticipada o solicitud satelital al único emisor}{04-documento-sin-cobertura}

Para SBD de la familia 9603, la documentación de Iridium permite hasta 340 bytes de subida y 270 de bajada por mensaje. Es capacidad por mensaje, no tamaño garantizado de un documento completo (Iridium, RockBLOCK 9603 \parencite{d3_sbd}). El protocolo propuesto usa solicitud correlacionada, versión, longitud, número de fragmento, integridad y reensamblado verificable. El original XML firmado se conserva en el contable y en su custodia; la cabina debe recibir todos los datos y el timbre necesarios para portar la representación válida aplicable, o el archivo de representación completo. No se exige transmitir todo el XML si no es necesario para esa representación; tampoco se acepta una representación incompleta generada desde folio/hash. El gateway satelital requiere un adaptador autorizado hacia la plataforma; no se presupone que Iridium hable directamente el protocolo de IoT Hub.

El límite total de emisión es 90 s (\caso{RT-09.01}{32}). Debe medirse la solicitud, emisión, retorno fragmentado, verificación y disponibilidad portable con el documento de mayor tamaño previsto. Para tamaño comprimido B y cabecera de h bytes por mensaje, con 0 $\leq$h <270, la bajada requiere al menos \texttt{ceil(B / (270 − h))} mensajes; el ensayo debe incluir pérdidas, reintentos y tiempo satelital. No se supone que un timbre completo quepa en una única trama.

La contingencia de falta de cobertura descrita por el SII no autoriza a audIT a emitir desde la cabina: la Restricción 8 conserva al contable como emisor. La representación y el envío fiscal deben seguir el régimen aplicable del contribuyente (SII, contingencia de guía electrónica \parencite{d3_sii_contingencia}; Resolución Exenta 99/2019 \parencite{d3_sii99}). La factibilidad del retorno SBD completo en 90 s queda como puerta de validación previa a aceptar ese camino; si no se logra, se revisan ADR-06 y T-11 con S3/D4. La regla inmediata es bloquear la salida sin documento conforme, sin afirmar cero detenciones.

La \figref{04-4-4-documento} detalla el intercambio y el cotejo previo al movimiento.

\anexoGrafico{../../Informe-2/D3/graficos/4-4-documento.png}{Secuencia de emisión, entrega y comprobación previa a la salida}{04-4-4-documento}

La salida sólo se habilita con carga coincidente, documento válido y representación portable disponible. Un cambio posterior de carga exige corregir mediante el emisor y volver a evaluar; la cabina no modifica la copia firmada ni reutiliza un folio para otro viaje.

\subsection{Fuentes con desfase y telemetría del vehículo}\label{sec:4-desfase}\label{sec:04-fuentes-con-desfase-y-telemetria-del-vehiculo}

Dos integraciones no entregan datos en el momento en que ocurren y la arquitectura las trata como tales. El consumo de combustible llega con hasta 40 días de desfase (\caso{numeral 7.3}{15}), y los peajes se liquidan mensualmente. Ambas se ingieren por lotes desde un depósito seguro, con validación sintáctica y verificación de totales, y ambas quedan detrás de un adaptador, de modo que el día en que una de esas contrapartes publique una interfaz en línea baste conectar el adaptador nuevo sin tocar el modelo de costeo.

Ingerir el archivo no basta. Cada carga de combustible se cruza contra la posición del camión en ese instante y cada pasada de peaje contra la traza del viaje. Ese cruce es lo que convierte un archivo de cobros en costo imputable a un viaje, y es también lo que deja a la vista la carga que se registró en una estación por la que el camión no pasó. Son 74.000 abastecimientos y 620.000 pasadas de peaje al año según el \caso{numeral 14.1}{29}, volumen que no admite revisión manual.

Un tercer frente es la posición de la flota de terceros. El Caso declara 340 de 374 camiones con dispositivo repartidos en tres plataformas distintas (\caso{numeral 14.1}{29}), dos de ellas con acceso de solo consulta y una que ni siquiera permite exportar (\caso{numeral 16.1}{34}). Unificar esa vista es obligación de esta oferta y reemplazar esos equipos está excluido (\caso{capítulo 11}{24}). Lo que se entrega es la vista unificada y el estándar de homologación contra el cual clasificar cada equipo, no un inventario de un parque que el Caso no describe.

La telemetría de fábrica de los 61 tractocamiones se lee por acoplamiento sin contacto sobre la interfaz del vehículo, en modo de solo lectura y sujeta a la autorización de cada fabricante que exige RT-17.06 del Caso. La ausencia de contacto eléctrico no demuestra por sí sola autorización ni disponibilidad de señales. Cada lectura requiere protocolo, PGN/parámetros, licencia de acceso y homologación del fabricante. Esa elección responde a las restricciones del Caso. La restricción 6 prohíbe que el equipamiento a bordo afecte la garantía del vehículo o interfiera con sus sistemas de seguridad, y el Capítulo 11 excluye intervenir la electrónica de fábrica. No cortar ni empalmar el arnés reduce la intervención, pero la conformidad de garantía y seguridad se verifica por modelo de vehículo; los 182 lectores previstos no implican 182 fuentes CAN disponibles. Los parámetros que se leen se listan en la \tabref{04-tabla-4-14}.

\begin{tabla}[ancho=completo, fuente=condensada]{Parámetros leídos de la telemetría de fábrica}{04-tabla-4-14}{L{47mm} X}
Parámetro & Para qué se usa \\ \midrule
Consumo instantáneo y acumulado & Componente medido del costo por kilómetro de la flota propia \\
Nivel de estanque & Contraste con el abastecimiento facturado por la red de estaciones \\
Revoluciones y aceleraciones bruscas & Conducción eficiente y explicación de la dispersión de rendimiento de 19 por ciento entre camiones del mismo modelo y ruta \\
Odómetro del vehículo & Kilómetro efectivo del viaje, denominador del costo \\
Horas de funcionamiento y de ralentí & Consumo que no produce kilómetro \\
Uso de freno de servicio y freno motor & Insumo del mantenimiento por condición \\
\end{tabla}

\notaTabla{Fuente: elaboración propia sobre el Caso, criterio 18.}

Todos son parámetros de lectura. Ninguno escribe sobre el bus del vehículo, y esa restricción es de diseño y no de configuración.

\subsection{Capa analítica y costo real por kilómetro}\label{sec:4-analitica}\label{sec:04-capa-analitica-y-costo-real-por-kilometro}

La segregación entre lo transaccional y lo analítico se resuelve por captura de cambios hacia un repositorio organizado en capas sucesivas de refinamiento, desde el dato crudo hasta el modelo dimensional que consume la gerencia. La replicación lee la bitácora del motor transaccional y no consulta sus tablas, para evitar consultas analíticas sobre OLTP. La lectura de WAL y sus slots consumen recursos y se monitorizan; la separación se comprueba bajo carga. Las capas del repositorio se describen en la \tabref{04-tabla-4-15}.

\begin{tabla}[ancho=completo, fuente=condensada]{Capas del repositorio analítico}{04-tabla-4-15}{L{36mm} X X}
Capa & Qué contiene & Qué garantiza \\ \midrule
Cruda & El registro tal como llegó, replicación transaccional, posición telemática y archivos mensuales de combustible y peaje & Reproceso completo sin volver a la fuente y trazabilidad del origen de cada indicador \\
Depurada & Deduplicación, validación de coordenadas, corrección de marcas de tiempo y cruce de la telemetría con el tramo de la orden & Un mismo hecho con una sola representación \\
De negocio & Modelo dimensional con el hecho de costo por viaje y por tramo, y las dimensiones de ruta, cliente, tracto, conductor, régimen de propiedad y tiempo & Consulta de gerencia sin conocimiento del modelo transaccional \\
\end{tabla}

\notaTabla{Fuente: elaboración propia.}

La \figref{04-audit-figura-4-5-analitica} muestra la capa analítica y la explotación del costo por kilómetro.

\anexoGrafico{../../Informe-2/D3/graficos/audit-figura-4.5-analitica.png}{Capa analítica por niveles de refinamiento y explotación del costo por kilómetro}{04-audit-figura-4-5-analitica}

Fuente: D4, figura conservada.

El costo por kilómetro se calcula distinto según el régimen de propiedad del equipo. Para la flota propia se compone de combustible medido por telemetría, peajes efectivamente transitados, neumáticos, mantenimiento, jornada del conductor y depreciación. Para la flota subcontratada el mandante solo conoce la tarifa pactada y los anticipos de combustible, de modo que el resto se imputa. Esa diferencia hay que declararla, porque comparar flota propia con flota de terceros sin declararla es comparar cosas distintas, y esa comparación gobierna la decisión de crecer con una u otra. La \tabref{04-tabla-4-16} muestra el origen de cada componente.

\begin{tabla}[ancho=completo, fuente=condensada]{Componentes del costo por viaje y su origen}{04-tabla-4-16}{L{36mm} X X}
Componente & Flota propia & Flota subcontratada \\ \midrule
Combustible & Medido por telemetría y valorizado con la liquidación mensual & No observable. La compañía solo conoce el anticipo que otorga \\
Peajes & Pasada efectiva cruzada con la traza & Igual, cuando el peaje lo asume la compañía \\
Conductor & Jornada imputada al tramo & Incluida en la tarifa pactada, no descomponible \\
Mantenimiento y neumáticos & Cuota por kilómetro sobre la orden de taller y la posición de neumático & No observable \\
Sobreestadía & Horas de espera acreditadas por geocerca & Igual \\
Tarifa a terceros & No aplica & Costo directo contractual \\
Denominador & Kilómetro del odómetro telemático & Kilómetro de la traza de posición \\
\end{tabla}

\notaTabla{Fuente: elaboración propia.}

El mecanismo es dual y no es una elección de diseño, es lo que el propio Caso fija. \caso{RT-05.29}{32} exige el costo consolidado de un viaje en no más de 24 horas tras su cierre, con los componentes que a esa fecha estén disponibles y con indicación explícita de los que aún no lo están. Se emite entonces una versión preliminar dentro de esas 24 horas, marcada como tal y enumerando qué falta, y una versión definitiva cuando llegan el combustible y los peajes. La primera no se sobrescribe. Ambas coexisten, y la desviación entre una y otra mide la calidad de la estimación. Sin ese doble paso el costo por viaje se convierte en un cierre contable tardío, que es exactamente la condición que permitió servir tres contratos bajo costo, el peor durante cuatro años.

El mismo RT-05.29 fija el resto de las latencias analíticas, y conviene tenerlas juntas porque gobiernan el diseño de la capa. La \tabref{04-tabla-4-17} las reúne.

\begin{tabla}[ancho=completo, fuente=condensada]{Latencias de la capa analítica fijadas por RT-05.29 del Caso}{04-tabla-4-17}{L{47mm} X}
Dato & Latencia máxima \\ \midrule
Posición de un camión con cobertura & 2 minutos \\
Jornada acumulada de un conductor & Tiempo real, disponible en el momento de asignar \\
Tiempo de llegada y de salida en un punto de cliente & Registrado en el momento del evento \\
Costo consolidado de un viaje & 24 horas tras su cierre, con indicación de los componentes pendientes \\
Emisiones & Consolidación mensual \\
\end{tabla}

\notaTabla{Fuente: \caso{RT-05.29}{32}.}

La jornada acumulada es la fila que condiciona la arquitectura. Exigirla en tiempo real en el momento de asignar significa que no puede leerse del repositorio analítico, y por eso la evaluación utiliza estado autoritativo o una proyección cuya versión y actualidad se comprueban antes del commit. Una caché antigua no habilita al conductor.

\subsection{Explotación analítica y autoservicio}\label{sec:4-autoservicio}\label{sec:04-explotacion-analitica-y-autoservicio}

RT-05.25 obliga a proveer tableros operacionales y de gestión sobre los indicadores que el Caso define. RT-05.26 exige poder filtrar por período y por dimensión propia del caso y profundizar desde el indicador agregado hasta la transacción de origen. Esa navegación termina en el viaje individual con su traza, su documento de transporte y su respaldo de entrega, y no en un subtotal.

RT-05.27 obliga a que el cliente construya sus propios informes sin intervención del adjudicatario, mediante una herramienta de autoservicio con modelo semántico documentado. El modelo semántico se entrega validado con la gerencia de administración y finanzas, con las métricas nombradas y definidas una sola vez, de modo que margen por ruta signifique lo mismo en todos los tableros. RT-05.28 exige que todo informe sea exportable en formatos abiertos y programable para envío automático por calendario.

RT-05.30 valora la analítica predictiva con el modelo, sus variables, su métrica de desempeño y su plan de reentrenamiento documentados. Es un requisito deseable y esta oferta lo aborda en el Subdocumento 13.

\subsection{Secuencia de liquidación y costeo por viaje}\label{sec:4-secuencia-costeo}\label{sec:04-secuencia-de-liquidacion-y-costeo-por-viaje}

El proceso de consolidación de costos operacionales y liquidación formal a transportistas se representa en la \figref{04-4-6-liquidacion}.

\anexoGrafico{../../Informe-2/D3/graficos/4-6-liquidacion.png}{Secuencia de estimación de costo por viaje, conciliación de insumos y liquidación ERP}{04-4-6-liquidacion}

Fuente: elaboración de audIT.

Como ilustra la \figref{04-4-6-liquidacion}, al finalizar el viaje el contexto de Liquidación emite la versión preliminar v1 en menos de 24 horas (\caso{RT-05.29}{31--32}), incorporando versiones incrementales auditadas (v2 a vn) a medida que ingresan las cartolas de combustible y telepeaje TAG, hasta consolidar la liquidación mensual definitiva en el sistema contable.

\subsection{Modelo táctico del dominio}\label{sec:4-tactico}\label{sec:04-modelo-tactico-del-dominio}

RT-02.13 exige presentar el modelo de dominio del negocio con las entidades principales, sus relaciones y los eventos de negocio que las modifican. El modelo táctico de la \figref{04-audit-figura-4-6-tactico} lo entrega organizado por contexto delimitado, de modo que cada agregado quede bajo el contexto que lo posee.

\anexoGrafico{../../Informe-2/D3/graficos/audit-figura-4.6-tactico.png}{Modelo táctico del dominio. Agregados, entidades y servicios por contexto}{04-audit-figura-4-6-tactico}

Fuente: D4, figura conservada.

Cada agregado de la figura pertenece a un solo contexto delimitado.

\subsection{Responsabilidad y emplazamiento de componentes lógicos}\label{sec:04-responsabilidad-y-emplazamiento-de-componentes-logicos}

La \tabref{04-tabla-4-18} sintetiza responsabilidades y emplazamiento; las cantidades y configuraciones exhaustivas corresponden a T-11. Cada contexto conserva sus datos bajo un propietario y expone contratos, sin acceder a las tablas de otro (\art{16.2}{11}).

\begin{tabla}[ancho=completo, fuente=condensada]{Componentes lógicos y fundamento del emplazamiento}{04-tabla-4-18}{L{36mm} X X}
Componente & Responsabilidad & Emplazamiento y motivo \\ \midrule
Planificación y tráfico & Orden, asignación y retornos & Servicios en AKS; requiere estado coherente de todos los recursos y plazo de 30 s \\
Flota y activos & Aptitud, semirremolque y taller & AKS; integra maestros acreditados y lectura autorizada. La captura instrumental reside a bordo \\
Personas y cumplimiento & Evidencia, jornada, habilitaciones y permisos & AKS más registro local; decide con identidad y evidencia conservadas, no con BI \\
Telemetría y geocercas & Captura, posición, llegada/salida y reconexión & Captura/evaluación de geocerca a bordo; ingesta y vista unificada en nube; funciona sin equipos en recintos del cliente \\
Operación de fletes & Viaje ocurrido, documento, conformidad y carga peligrosa & AKS más aplicación de terreno; el emisor tributario permanece externo detrás de ACL \\
Liquidación y costeo & Versiones de costo y liquidación & Servicio y repositorio analítico separado; contable recibe sólo registros aprobados y traducidos \\
PostgreSQL / TimescaleDB / Redis & Negocio, series y caché, respectivamente & Instancias separadas conforme a S5 y T-11; Redis no es autoridad única de despacho \\
IoT Hub / Event Hubs / outbox & Admisión, eventos y publicación duradera & Nube con identidad por dispositivo, cuotas, partición, reintento y deduplicación \\
Búfer, móvil y terminal & Captura sin enlace y copia documental & Equipo local y aplicación por perfil; autonomía y conflictos descritos en 4.1.26 \\
\end{tabla}

La \tabref{04-tabla-4-18} usa los seis nombres canónicos de S3. Los objetivos de latencia son los del Caso; estimaciones de milisegundos de un componente se evalúan como hipótesis y no como prueba de cumplimiento extremo a extremo.

\subsection{Concurrencia y volumen declarados}\label{sec:4-concurrencia}\label{sec:04-concurrencia-y-volumen-declarados}

\caso{RT-09.02}{32} no fija un número. Ordena derivarlo de la volumetría del numeral 14.1 y declararlo conforme al numeral 14.2, considerando de manera expresa la reconexión simultánea de unidades al salir de zonas de sombra. La derivación de la concurrencia está en la \tabref{04-tabla-4-19}.

\begin{tabla}[ancho=completo, fuente=condensada]{Concurrencia derivada de la volumetría del Caso}{04-tabla-4-19}{L{36mm} X X}
Población & En hora punta & Base de la derivación \\ \midrule
Personal interno con acceso a sistemas & 50 a 80 & Torre de 22 personas en turno, despacho de terminal, finanzas, facturación y taller, sobre los 336 con acceso del numeral 14.1 \\
Conductores en aplicación móvil & 100 a 150 & Accesos breves de inicio y cierre de turno sobre 454 conductores \\
Transportistas subcontratados en portal & 30 a 50 & Consulta de viajes y de liquidación sobre 148 \\
Clientes en seguimiento & 50 a 100 & Sesiones de seguimiento sobre 84 clientes activos \\
Suma aritmética de los rangos & 230 a 380 & Extremo inferior y superior de las cuatro filas \\
Concurrencia de dimensionamiento preliminar & 380 & Hipótesis de D3/D1 obtenida al sumar los extremos superiores de los cuatro rangos; no es una medición del Caso ni una concurrencia contractual. S4.2 de D4 mantiene 350 (525 en estrés), mientras D1/T-17 utiliza 380 (570); D3 y D4 deben acordar una única línea base y comunicarla a S7/S9 antes de dimensionar o declarar aceptación. \\
\end{tabla}

\notaTabla{Fuente: elaboración propia sobre el \caso{numeral 14.1}{29}.}

El volumen anual y los perfiles de ráfaga se distinguen en la \tabref{04-tabla-4-20}. Son supuestos de capacidad solicitados por \caso{numeral 14.2}{30}; se reemplazan por mediciones en el levantamiento.

\begin{tabla}[ancho=completo, fuente=condensada]{Derivación reproducible de la carga de datos}{04-tabla-4-20}{L{36mm} X X}
Magnitud & Fórmula y resultado & Alcance del supuesto \\ \midrule
Horas anuales en marcha & 41.000.000 km / 55 km/h = 745.455 h & Velocidad comercial de diseño; el kilometraje es dato de FEP03 14.1 \\
Posición en marcha & 745.455 h $\times$ 120 muestras/h $\approx$ 89,45 millones & Cada 30 s; escenario de vista completa futura de 374 unidades, no acceso actual garantizado \\
Posición detenida & (374 $\times$ 8.760 − 745.455) h $\times$ 12 muestras/h $\approx$ 30,37 millones & Cada 5 min; año de 365 días y unidades dentro del alcance autorizado \\
Posición total y bytes crudos & 119,82 millones $\times$ 64 bytes $\approx$ 7,67 GB/año & Registro binario supuesto; JSON, índices, transporte y réplicas se miden aparte \\
Motor, escenario inicial de 61 unidades & (745.455 $\times$ 61 / 374) h $\times$ 60 $\times$ 160 bytes $\approx$ 1,17 GB/año & Reparto proporcional de marcha; autorización y señales reales por fabricante. No se asume CAN activo en toda la flota \\
Capacidad normal máxima de admisión & 374/30 + 182/60 $\approx$ 15,5 eventos/s, más eventos de negocio & Escenario simultáneo de diseño; no TPS de despacho ni cantidad de fuentes acreditadas \\
Asignación en punta & 22 puestos $\times$ 1 asignación/30 s $\approx$ 0,73 TPS; ráfaga propuesta 10$\times$ $\approx$ 7,33 TPS & Hipótesis de carga; no se deriva dividiendo viajes anuales por usuarios \\
Reconexión de 300 unidades, perfil 72 h de S4/T-17 & 1.786.200 registros y 971.296.800 bytes con fotos; $\geq$1.861 registros/s en 960 s útiles & 20 min menos 20 \% de reserva; ancho útil con overhead/reserva $\approx$10,12 Mbit/s; prueba CP-PERF-02 \\
Retención de series en línea & 2 años; posición cruda $\approx$15,34 GB y motor inicial $\approx$2,34 GB & Dimensionar además índices, WAL, réplicas, fotos, documentos y crecimiento. Factor físico provisional 3$\times$, sujeto a medición \\
\end{tabla}

La \tabref{04-tabla-4-20} no equivale al tamaño de toda la base de datos: incluye supuestos de muestreo, población y tamaño de registro. Los volúmenes transaccionales e históricos se estiman en el anexo de S5. El perfil por camión de 72 h y la reserva de búfer propuesta de 288 h se conservan separados del volumen anual. El cierre vial de doce días no se equipara automáticamente a doce días sin enlace.

\caso{RT-05.10}{31} exige dos años en línea para posición/telemetría, con agregación declarada para el resto. La carga de crecimiento a tres veces la volumetría inicial se prueba según RT-09.03 transversal, además de estrés a 570 sesiones y ráfaga de 300 reconexiones; son ensayos distintos.

\subsection{Modelado de amenazas de la arquitectura}\label{sec:04-modelado-de-amenazas-de-la-arquitectura}

Para la baliza Bluetooth asociada al semirremolque considerada en la innovación, se modelan a nivel de diseño las amenazas STRIDE de la \tabref{04-tabla-4-21}. Este análisis responde a \transv{RT-26.07}{44}; define controles propuestos y no constituye evidencia de pruebas de seguridad ejecutadas. Si cambia el protocolo o el alcance de la baliza, el modelo debe actualizarse antes de su aprobación.

\begin{tabla}[ancho=completo, fuente=condensada]{Amenazas y controles de diseño para la baliza y su integración}{04-tabla-4-21}{L{36mm} X X}
Categoría STRIDE & Amenaza considerada & Control de diseño y efecto esperado \\ \midrule
Suplantación & Baliza desconocida o asociada a otro semirremolque/viaje & Vincular el identificador registrado con el activo y el viaje vigente; marcar eventos desconocidos o fuera de contexto y excluirlos de la evidencia habilitante \\
Manipulación & Evento alterado, repetido o fuera de secuencia & Validar esquema y vigencia temporal, conservar procedencia, aplicar deduplicación e idempotencia; comprobar secuencia solo si el protocolo entrega ese dato \\
Repudio & Disputa sobre el origen o tratamiento del evento & Registrar identificador de fuente, hora de recepción, correlación con viaje y acciones posteriores; enlazar las rectificaciones al registro original. La señal de baliza por sí sola no acredita jornada ni descanso legal \\
Divulgación de información & Exposición de ubicación o asociación entre conductor y semirremolque & Limitar la consulta a roles autorizados y al viaje/cliente habilitado; minimizar los datos expuestos y proteger su transmisión y almacenamiento en la plataforma \\
Denegación de servicio & Pérdida, retraso o ráfaga de eventos de baliza & Controlar tasa y vigencia, señalar explícitamente la ausencia o antigüedad de la última señal y degradar la vista sin inventar posiciones ni bloquear otras fuentes disponibles \\
Elevación de privilegios & Evento de baliza usado para alterar despacho o controlar el vehículo & Mantener la ingesta como entrada de solo lectura; no otorgar a la baliza facultad de actuar sobre CAN ni de cambiar por sí sola estados del viaje o decisiones de despacho \\
\end{tabla}

La \tabref{04-tabla-4-21} acota la baliza como fuente auxiliar: la operación conserva la procedencia del evento y la evaluación independiente de las reglas de despacho y jornada.

\subsection{Estilo, principios de diseño y contratos de dominio}\label{sec:04-estilo-principios-de-diseno-y-contratos-de-dominio}

El estilo combina servicios independientes, transacciones locales y eventos. Un monolito modular ofrece menor costo de operación, pero no satisface el despliegue independiente exigido para los componentes críticos; una cadena síncrona de servicios multiplica tiempos de espera; SOA con orquestación central concentra fallas. Los seis contextos con API para decisiones inmediatas y eventos para efectos derivados separan la carga de \textasciitilde{}15,5 eventos/s normal de la ráfaga de reconexión de \textasciitilde{}1.861 registros/s, sin deducir el estilo únicamente del volumen anual.

Se aplica SOLID de forma verificable: cada servicio mantiene una responsabilidad de contexto (SRP); nuevos proveedores se incorporan mediante adaptadores del mismo contrato (OCP); todo adaptador debe superar las pruebas comunes de ese contrato (LSP); las interfaces de jornada, permisos y liquidación se segregan según su consumidor (ISP); el dominio depende de puertos de repositorio/mensajería y no del SDK cloud (DIP). Se documentan vistas y preocupaciones bajo ISO/IEC/IEEE 42010 y gobierno arquitectónico equivalente a TOGAF mediante ADR, Comité de Arquitectura y revisión mensual (\art{4.3}{5}).

La \tabref{04-tabla-4-22} completa la \figref{04-audit-figura-4-6-tactico} con comandos y eventos. Los nombres de contratos se proponen para implementación; cada evento lleva ID, agregado, versión, secuencia y correlación, con esquemas AsyncAPI y política de compatibilidad.

\begin{tabla}[ancho=completo, fuente=condensada]{Responsabilidad, comandos y eventos por contexto}{04-tabla-4-22}{L{36mm} X X}
Contexto & Comando principal & Evento derivado y consumidor \\ \midrule
Planificación y tráfico & AsignarRecursos / ReprogramarViaje & ViajeAsignado $\rightarrow$ Operación, Telemetría; rechazo conservado en auditoría \\
Flota y activos & AcreditarVigencia / RegistrarIntervencion & AptitudEquipoActualizada $\rightarrow$ Planificación; invalida proyección afectada \\
Personas y cumplimiento & EvaluarAsignacion / RevocarPermiso & JornadaActualizada y PermisoRevocado $\rightarrow$ Planificación y canales autorizados \\
Telemetría y geocercas & IngerirLote / EvaluarGeocerca & LlegadaRegistrada / SalidaRegistrada $\rightarrow$ Operación; posición autorizada al portal \\
Operación de fletes & SolicitarDocumento / CerrarViaje & DocumentoDisponible / ViajeCerrado $\rightarrow$ Planificación y Liquidación \\
Liquidación y costeo & CalcularCosto / AprobarLiquidacion & CostoVersionado / LiquidacionAprobada $\rightarrow$ portal y ACL contable \\
\end{tabla}

La \tabref{04-tabla-4-22} no requiere una transacción distribuida entre seis bases. Planificación valida versiones y reservas en su transacción; los contextos conservan sus agregados y publican efectos desde outbox. Cambios invalidantes concurrentes fuerzan una nueva evaluación antes de confirmar.

\subsection{Trazabilidad del alcance de Informe 2}\label{sec:04-trazabilidad-del-alcance-de-informe-2}

El T-12 vigente mantiene los 28 RF y 14 RNF, sus paquetes EDT y casos CP asociados. Los contratos de S4.1 se corresponden con ese catálogo. Los componentes se ubican en la \tabref{04-tabla-4-18} y la \figref{04-logicacapas}. RF-001 atraviesa Planificación, Personas y Flota; RF-013/RF-014 terminan en ACL y contable; RNF-008 atraviesa borde, ingesta y persistencia. Esta relación permite comprobar cobertura del alcance sin convertir el cuerpo en inventario de formularios.

\subsection{Seguridad lógica y superficie de exposición}\label{sec:04-seguridad-logica-y-superficie-de-exposicion}

El acceso público termina en Front Door/WAF y API Management; bases, bus, secretos y servicios internos permanecen en red privada conforme a RT-03.04. El seguimiento sin autenticación muestra sólo el estado del envío mediante número de documento, sin posición personal ni historial del transportista (\caso{RT-16.30}{33}). Los portales autenticados evalúan rol, transportista, recurso, cliente, viaje activo, finalidad y permiso vigente en cada petición.

Entra ID autentica personas y los equipos usan identidad/certificado individual revocable. El conductor se identifica estando detenido, mediante tarjeta/token asignado y vinculación a persona y viaje; la validación con guantes y rotación de conductores pertenece a S9, sin exigir una contraseña recordada ni interacción en marcha (RT-12.11/RT-13.08 del Caso). Identidades administradas acceden a Key Vault; claves y secretos no aparecen en código ni URI. Se cifra el transporte y los campos exigidos por RT-11.10, con separación de claves y roles. Los accesos sensibles se registran con qué se mostró y autorización utilizada, conforme a RT-16.09 del Caso.

Se verifican, como mínimo, suplantación de dispositivo y baliza, alteración/repetición de eventos, repudio de firma, consulta entre transportistas, saturación de reconexiones y elevación a administración. La \tabref{04-tabla-4-21} cubre STRIDE de la baliza; las pruebas del resto se trazan a S9/T-17. Ni el hash ni el cifrado por sí solos acreditan cumplimiento legal o seguridad integral.

\subsection{Ambientes y liberación}\label{sec:04-ambientes-y-liberacion}

DEV usa datos sintéticos; QA contiene escenarios y simuladores reproducibles; PREPROD reproduce integraciones y configuración sin accesos de producción; PROD contiene los datos autorizados y cambios aprobados; DR replica los componentes críticos bajo el mecanismo jurídico y de continuidad de S5/S4.3. Cada ambiente tiene identidades, secretos, redes, almacenes y permisos separados. La configuración se versiona como código y se promueve el mismo artefacto firmado según S6.

La aceptación incluye pruebas de contratos, carga, reconexión, evidencia, permisos y reversión. Activar Brazil South con datos personales exige un mecanismo de transferencia documentado; deshabilitar esa réplica sin otro sitio secundario operativo tampoco acredita RPO $\leq$15 min/RTO $\leq$4 h. La revisión jurídica y la conmutación probada son puertas conjuntas de producción.

\subsection{Operación desconectada, reconexión y escalamiento}\label{sec:04-operacion-desconectada-reconexion-y-escalamiento}

A bordo se conservan posición, eventos de jornada, llegada/salida y documentos durante al menos 72 h. En terminal se adopta el mayor mínimo: 24 h del RT-03.10 transversal frente a 12 h del Caso. La captura local emplea SQLite/WAL y outbox; la aplicación comparte identidad y catálogos vigentes con el nodo autorizado. La operación local no inventa vigencias, permisos ni asignaciones globales. Toda nueva asignación requiere acceso al servicio autoritativo de reserva. RN-04 sólo excepciona una fuente de evidencia indisponible bajo sus controles; no sustituye la reserva, la legalidad ni el documento conforme necesario para el movimiento.

La \figref{04-4-7-reconexion} muestra la reconexión: autenticación, lotes numerados, admisión, deduplicación, commit y acuse duradero. Sólo el último acuse permite depurar el búfer. La conciliación conserva el original y la resolución de conflictos; el timestamp de recepción no reemplaza la hora del evento.

\anexoGrafico{../../Informe-2/D3/graficos/4-7-reconexion.png}{Reconexión con acuse posterior a persistencia y conciliación}{04-4-7-reconexion}

CP-PERF-02 debe demostrar $\leq$20 min para cada camión después de 72 h, con 300 unidades simultáneas y sin pérdida de jornada ni tiempos de cliente. La admisión por IoT Hub no demuestra entrega final. Las 288 h propuestas y la capacidad de búfer de S4.2 requieren otro ensayo; no sustituyen esta prueba contractual.

El escalamiento propuesto usa mínimo dos réplicas por contexto, aumento al superar CPU 65 \% o retraso de consumo 30 s durante dos minutos, con reducción sólo tras diez minutos estables. Máximo inicial: ocho réplicas de ingesta y cuatro por servicio de negocio, sujeto al perfil y cuotas de S4.2/T-11. Los mamparos reservan recursos para despacho, documento y emergencia. Al alcanzar el máximo se alerta a operación y se limita carga no crítica; no se incrementa capacidad ni costo sin registro. El impacto económico de esos límites corresponde al Sobre 3. S9 mide latencia completa, saturación, costo de recursos y recuperación antes de aceptar el perfil.

\section{Arquitectura física}\label{sec:4-fisica}
% ============================================================================

El \art{16.1}{11} exige una solución híbrida: la carga principal en nube pública y componentes
desplegados en las instalaciones u operaciones del mandante. Esta solución reparte sus componentes en
los tres planos de la \tabref{4-planos}. Cada componente lógico de la sección~4.1 queda en uno de ellos
por una razón que se declara en la sección~4.2.3.

\begin{tabla}{Los tres planos de emplazamiento}{4-planos}{L{32mm} X X}
Plano & Contenido & Por qué no puede estar en otro lugar \\ \midrule
Nube & Núcleo transaccional, analítica, portales, ingesta e integración, en Azure Chile Central con recuperación en Brazil South & Elasticidad hacia 430 camiones y absorción del peak de reconexión \\
On-premise de sitio & Nodo de continuidad de la torre en San Bernardo y gabinetes en los cuatro terminales regionales & RT-06.01 del Caso exige gabinete por terminal y el Art. 16.4 pide operar al menos 24 horas sin enlace \\
On-premise distribuido & 182 equipos audIT a bordo, más la lectura de los equipos que ya tienen 192 camiones de terceros & La operación no puede depender de la cobertura móvil (restricción 4) \\
\end{tabla}
\notaTabla{Fuente: FEP01, Artículo 16, pp. 11 y 12, y Caso, RT-06.01, p. 32, y capítulo 10, p. 23.}

El reparto de la flota sale del \caso{numeral 2.2}{6}: 340 de los 374 camiones tienen dispositivo y
«34 camiones subcontratados no tienen ninguno». Entonces 226 menos 34 da 192 camiones de terceros con
equipo, y 340 menos 192 da 148 propios, todos equipados. Los 148 propios y los 34 de terceros sin
equipo llevan el equipo de audIT, 182 en total. Los 192 de terceros que ya tienen equipo lo conservan,
porque la restricción 3 impide intervenirlo sin acuerdo de su dueño \caso*{capítulo 10}{23}, y sus
datos llegan por las plataformas de sus proveedores. La \figref{4-fisicafig} muestra los tres planos y
los enlaces que los unen.

\figuraHorizontal{figuras/04-arquitectura/fisica-general.png}{Vista general de la arquitectura física: nube primaria y secundaria, San Bernardo, terminales regionales y flota, con los enlaces y su capacidad}{4-fisicafig}

La figura se lee de arriba hacia abajo. Arriba está la nube, con Azure Chile Central como región
primaria en tres zonas y Brazil South en espera. Al medio está San Bernardo, conectado a la red central
de Azure por ExpressRoute de 100~Mbit/s y por una VPN de respaldo a través de otro proveedor. A su
lado están los cuatro terminales regionales, cada uno con un enlace fijo de al menos 10~Mbit/s y un
router 5G de respaldo. Abajo está la flota: los 182 equipos audIT transmiten por la red celular y,
donde no hay señal, por la red satelital Iridium, mientras que los 192 equipos de terceros entregan su
posición a través de las plataformas de sus proveedores. Las secciones que siguen desarrollan cada
parte con su propio diagrama, como pide el \comunicado{10}{sección 4.4}.

% ----------------------------------------------------------------------------
\subsection{Especificaciones Implementos a proveer (Hardware y Software)}\label{sec:4-implementos}
% ----------------------------------------------------------------------------

El hardware lo adquiere el mandante, y audIT debe especificar exactamente qué comprar, cuánto y con
qué características \caso*{capítulo 11}{24}. \transv{RT-08.10}{19} pide para cada dispositivo marca,
modelo de referencia, cantidad, características mínimas y costo unitario estimado. El costo va en el
Sobre N.º~3, porque el \art{50.2}{29} no admite en la Oferta Técnica cifras que permitan inferir el
monto. La \tabref{4-implementos} resume los implementos por sitio. El detalle, con características,
ubicación, ciclo de vida y justificación de cada uno, está en el Formulario T-11.

\begin{tabla}{Implementos que compra el mandante, por sitio}{4-implementos}{L{27mm} L{40mm} X L{24mm}}
Sitio & Implemento & Modelo de referencia & Cantidad \\ \midrule
Cabina & Equipo a bordo & iWave G26I con eMMC de 8 GB & 182 + 19 \\
Cabina & Módem satelital & Iridium Edge & 182 + 19 \\
Cabina & Lector CAN sin contacto & Technoton CANCrocodile & 182 + 19 \\
Cabina & Lector de tarjeta & GAO RFID MIFARE DESFire & 182 + 19 \\
Conductores & Tarjeta de identificación & MIFARE DESFire & 572 \\
Semirremolques & Baliza Bluetooth & Teltonika EYE Sensor BTSMP1 & 210 + 21 \\
Terminales & Computador de terminal & OnLogic Karbon 430 & 8 \\
Terminales & Router 5G de respaldo & Teltonika RUTX50 & 4 \\
Terminales & Switch de gabinete & Teltonika TSW202 & 4 \\
Terminales & UPS de gabinete & APC SRT1000XLI & 4 \\
Cinco terminales & Punto de acceso de patio & Fortinet FortiAP 234G & 5 \\
Cinco terminales & Lector de portería & Minew G1 & 5 \\
San Bernardo & Servidor & Dell PowerEdge R360 & 2 \\
San Bernardo & Firewall & Fortinet FortiGate 90G & 2 \\
San Bernardo & Switch & Fortinet FortiSwitch 124F-POE & 2 \\
San Bernardo & UPS & APC SRT3000RMXLI con SRT96RMBP & 2 \\
San Bernardo & Clima de precisión & Vertiv Liebert Mini-Mate2 MMD12E & 2 \\
San Bernardo & Detección y extinción & VESDA VLF-250 y Kidde Fluoro-K & 1 y 1 \\
San Bernardo & Control de acceso & Suprema FaceStation F2 & 2 \\
San Bernardo & Videovigilancia & Axis M3215-LVE y grabador S3008 Mk II & 4 y 1 \\
San Bernardo & Detección de agua & Vertiv Liebert LT460 & 1 \\
San Bernardo & Grupo electrógeno & Himoinsa HYW-8 T5 S5 con estanque de 100 L & 1 \\
\end{tabla}
\notaTabla{Fuente: elaboración propia. Detalle en el Formulario T-11 (AUDIT-Formulario-T-11.pdf). Los servicios de nube se declaran en la sección 4.2.4.}

Las cantidades salen de las poblaciones del Caso y no de la adhesión. El equipo a bordo y sus tres
periféricos se compran para 182 camiones, más 19 de reposición, que son el 10~\% del parque
redondeado hacia arriba. Las 572 tarjetas cubren a los 520 conductores que el \caso{numeral 14.1}{29}
proyecta a tres años, más el mismo 10~\%. Las balizas cubren los 210 semirremolques propios más 21 de
reposición. El T-11 declara además, en filas aparte, los equipos que la flota suma a tres años según el
mismo numeral: 22 equipos a bordo para los camiones propios que pasan de 148 a 170, y 35 balizas para
los semirremolques que pasan de 210 a 245. Los 192 equipos de terceros no figuran porque no se
reemplazan.

\paragraph{Estándar de homologación de los equipos de terceros} El Caso no pide reemplazar las
plataformas de posicionamiento instaladas en camiones de terceros, pero sí «especificar qué se
requeriría si hubiera que homologarlas» \caso*{capítulo 11}{24}. La \tabref{4-homologacion} fija ese
estándar mínimo. Contra él se clasifica cada equipo durante la Etapa~1, porque el Caso no entrega
marca, modelo ni capacidad de esos equipos y esta oferta no los supone.

\begin{tabla}{Estándar mínimo para homologar un equipo de terceros}{4-homologacion}{L{40mm} X L{36mm}}
Exigencia & Valor mínimo & Fuente \\ \midrule
Operación sin cobertura & 72 horas continuas sin pérdida de registros de posición y jornada & Caso, RT-03.10, p. 31 \\
Salida de datos & Interfaz programática con autenticación entre sistemas y sin clave estática en la dirección web & FEP02, RT-05.18, p. 12 \\
Frecuencia de posición & Cada 30 segundos en marcha y cada 5 minutos detenido & Supuestos de la sección 4.2.9 \\
Red móvil & LTE en las bandas que usan los operadores en Chile, sin depender de 2G ni de 3G & SUBTEL (2026) y Diario Financiero (2026) \\
Identificación del conductor & Sin manipular un dispositivo en marcha ni recordar una credencial & Caso, RT-12.11, p. 32 \\
Instalación & Sin cortar ni empalmar el arnés del vehículo & Restricción 6, Caso, p. 24 \\
\end{tabla}
\notaTabla{Fuente: elaboración propia sobre las bases citadas y \textcite{subtel_bandas}, \textcite{df_3g_entel}.}

Las plataformas de los proveedores preexistentes de los 192 camiones subcontratados se integran mediante conectores y adaptadores API REST/Webhook hacia la capa de ingesta cloud de audIT SpA. Para aquellas unidades cuyo equipamiento preexistente no cumpla con los requisitos mínimos de almacenamiento desconectado de 72 horas (RT-03.10) o presente obsolescencia por el apagado de redes 2G/3G, el plan de adhesión del Subdocumento~3 contempla una opción de actualización tecnológica (retrofit) mediante la entrega en comodato de unidades audIT homologadas durante el paso normal por terminal, asegurando la cobertura integral de la flota sin fricción patrimonial para los transportistas.

% ----------------------------------------------------------------------------
\subsection{Equipo a bordo y almacenamiento derivado}\label{sec:4-bordo}
% ----------------------------------------------------------------------------

\caso{RT-06.01}{32} ordena tratar el dispositivo a bordo como un componente on-premise distribuido,
con su propio ciclo de vida, su mecanismo de actualización remota, su gestión de seguridad y su plan
de reposición, todo ello sujeto a que solo puede intervenirse físicamente cuando el camión pasa por un
terminal. El equipo elegido es el iWave G26I, un computador vehicular con Linux embebido, tres puertos
CAN con J1939, RS232 y RS485, red celular LTE Cat~4 o Cat~M1, GNSS, Wi-Fi, Bluetooth~5.0 y el elemento
seguro Microchip TA100 para arranque seguro y almacenamiento de claves. Funciona de 9 a 32~V, de
$-40$ a $+70$~°C y con protección IP67 \parencite{iwave_g26i}, lo que responde a la vibración, el polvo
y la temperatura de cabina que piden \transv{RT-08.11}{19} y \transv{RT-08.12}{19}. La
\figref{4-bordofig} muestra el equipo con sus periféricos.

\figura{figuras/04-arquitectura/equipo-a-bordo.png}{Equipo a bordo: iWave G26I, periféricos, interfaces y particiones de la memoria}{4-bordofig}

Cada periférico de la figura responde a una restricción del Caso. El CANCrocodile lee el bus del camión
a través del aislamiento de los cables, sin cortarlos, y entrega la señal en CAN~2.0B según SAE~J1939
\parencite{technoton_cancrocodile}, con lo que se cumple la restricción~6 sobre la garantía del
vehículo \caso*{capítulo 10}{24}. El lector de tarjetas MIFARE DESFire, IP66 y conectado por RS485
\parencite{gao_desfire}, identifica al conductor cuando acerca su tarjeta con el camión detenido, como
pide \caso{RT-12.11}{32}. El Iridium Edge, conectado por RS232, manda mensajes de hasta 340 bytes por
la red satelital \parencite{groundcontrol_edge} en los tramos de más de 80~km sin señal de la
restricción~4 \caso*{capítulo 10}{23}. Es el mismo módem que Webfleet vende en Chile dentro de su
producto satelital \parencite{webfleet_sat_cl}, lo que hace factible su homologación local. Ese
producto guarda hasta 40 horas de registro y exige su propio rastreador, y por eso el módem se conecta
al G26I, que sí guarda 288 horas. El Wi-Fi del equipo solo se usa en el patio del terminal para las
actualizaciones, que se aplican en dos particiones de sistema para volver a la anterior si la nueva
falla \parencite{mender_variables}, y la restricción~5 queda cumplida porque nada se instala ni se
actualiza fuera de un terminal. El Bluetooth lee la baliza del semirremolque acoplado, que es la base
de la innovación tecnológica del Subdocumento~13. La \figref{4-instalacionfig} ubica cada pieza en el tractocamión y en el semirremolque.

\figuraHorizontal{figuras/04-arquitectura/instalacion-a-bordo.png}{Distribución física del equipamiento embarcado en el tractocamión y en el semirremolque}{4-instalacionfig}

Esta disposición física garantiza que los sensores y antenas operen con línea de vista y máxima cobertura sin alterar la garantía del fabricante del vehículo.

La compra exige que el modelo cubra las bandas LTE que usan los operadores en Chile, de 700, 850,
1.700, 1.900 y 2.600~MHz \parencite{subtel_bandas}, y que cuente con la homologación de SUBTEL. La ficha
del G26I declara certificaciones CE, FCC, ISED, PTCRB y AT\&T, y no declara E-Mark ni SUBTEL, de modo
que ambas quedan como condición de compra. El equipo a bordo no emite documentos tributarios: la
restricción~8 deja al sistema contable como único emisor \caso*{capítulo 10}{24}.

\paragraph{Capacidad de almacenamiento} La observación 73 exige derivar la memoria sin ajustar el cálculo a los 8 GB del producto. Se distingue volumen serializado, sobrecarga física, reserva y particiones del sistema. El perfil de 72 h se calcula en \verSeccionFolio{sec:4-dimensionamiento}; 288 h es ampliación propuesta por audIT, independiente de la duración del cierre fronterizo.

La carga serializada del perfil suma 3.237.656 bytes por 72 h; cuatro repeticiones suman 12.950.624 bytes. Se presupuesta otro tanto para índices, WAL, cifrado y metadatos, y otro tanto como reserva: 38.851.872 bytes. Son presupuestos de ingeniería, no ocupación medida ni compresión conseguida. Una partición de 64 MiB aporta 67.108.864 bytes y deja 28.256.992 bytes adicionales sobre ese presupuesto; HIL debe comprobar ocupación, reinicio e integridad del perfil completo. Si la sobrecarga real supera el presupuesto, se recalcula antes de homologar.

La \tabref{4-emmc} suma asignaciones físicas en MiB (\(1\,\mathrm{MiB}=1.048.576\text{ bytes}\)); GB comercial se expresa en decimal. Así se evita sumar MB y MiB como si fueran equivalentes.

\begin{tabla}{Asignación de almacenamiento a bordo por componente}{4-emmc}{X R{27mm} L{58mm}}
Componente & Asignación & Base de diseño \\ \midrule
Partición de arranque & 16 MiB & Gestor de actualización \\
Sistema A y sistema B & 2.048 MiB & Hasta 1 GiB por imagen \\
Búfer propuesto de 288 h & 64 MiB & 38.851.872 bytes presupuestados \\
Diagnóstico & 256 MiB & Tope rotativo independiente \\
Geocercas y maestros & 16 MiB & Catálogo y vigencias \\
Total asignado & 2.400 MiB & 2.516.582.400 bytes \\
\end{tabla}
\notaTabla{Fuente: elaboración propia; mínimo de 72 h en Caso RT-03.10, p. 31. Hipótesis y reserva verificables mediante T-17.}

La asignación suma aproximadamente 2,52 GB decimales. Los 8 GB nominales del G26I \parencite{iwave_g26i} son una configuración de producto; la capacidad útil después de formateo y las reservas del controlador se verifican antes de compra. No se toma la capacidad nominal como toda la memoria disponible. El presupuesto contiene las dos imágenes y no se limita a los bytes de telemetría. La reserva para almacenamiento no demuestra vida útil de escritura: se contrastan escrituras, amplificación y resistencia del fabricante durante los 56 meses. Los plazos legales de conservación corresponden al repositorio central por dominio; el búfer local no reemplaza esos plazos.

% ----------------------------------------------------------------------------
\subsection{Emplazamiento de componentes}\label{sec:4-emplazamiento}
% ----------------------------------------------------------------------------

El \art{16.2}{11} obliga a justificar el emplazamiento componente por componente en función de
latencia, criticidad operacional, volumen de datos, restricciones regulatorias, disponibilidad de
conectividad y costo total de propiedad, y califica como observación grave una asignación no
justificada. La justificación de cada componente, con esos seis criterios, está en el Formulario T-11.
La \tabref{4-reparto} resume dónde queda cada uno.

\begin{tabla}{Reparto de los componentes por emplazamiento}{4-reparto}{L{12mm} X L{40mm} R{16mm}}
Código & Emplazamiento & Tipología & Componentes \\ \midrule
N & Azure Chile Central, tres zonas & No aplica & 23 \\
N2 & Azure Brazil South, en espera & No aplica & 1 \\
SB & San Bernardo, sala de 26 m² & Sala técnica secundaria o de sitio & 4 \\
GT & Gabinete de cada terminal regional & Gabinete o borde operacional & 3 \\
DB & Equipo a bordo de 182 camiones & On-premise distribuido & 5 \\
SR & Baliza en 210 semirremolques & On-premise distribuido & 1 \\
BM & Teléfono del conductor & No aplica & 1 \\
\end{tabla}
\notaTabla{Fuente: elaboración propia. Detalle por componente en el Formulario T-11.}

De 38 componentes, 23 quedan en la región primaria de nube, uno en la secundaria y 14 fuera de la
nube, repartidos entre San Bernardo, los gabinetes, los camiones, los semirremolques y el teléfono del
conductor. La propuesta cumple el \art{16.1}{11}, porque no es exclusivamente en nube ni
exclusivamente on-premise, y la parte on-premise sostiene las tres funciones que \caso{RT-21.06}{34}
clasifica en severidad máxima cuando cae el enlace: asignar un viaje, emitir un documento de transporte
y recibir un evento de emergencia. El Informe~1 declaraba 18 componentes en nube y su tabla sumaba 20
(observación 71). La cuenta actual sale de la lista del T-11.

\paragraph{Trazabilidad entre la arquitectura lógica y la física} El \comunicado{10}{sección 11}
recomienda una matriz que cruce cada componente con su lugar en el esquema de solución, en la
arquitectura lógica y en el nodo físico. La \tabref{4-trazabilidad} la presenta para los componentes
críticos.

\begin{tabla}{Trazabilidad de los componentes críticos entre la sección 4.1 y el nodo físico}{4-trazabilidad}{L{40mm} L{30mm} X L{16mm}}
Componente & Capa lógica (4.1) & Nodo físico (4.2) & T-11 \\ \midrule
Despacho y asignación & Servicios de negocio & AKS en tres zonas de Chile Central, con copia local en el nodo de San Bernardo & N y SB \\
Personas y cumplimiento & Servicios de negocio & AKS en tres zonas de Chile Central & N \\
Bus de telemetría & Integración y eventos & IoT Hub S1 y Event Hubs Premium en la subred de integración & N \\
Motor transaccional & Datos & PostgreSQL Flexible con alta disponibilidad entre zonas & N \\
Series de tiempo & Datos & PostgreSQL Flexible con TimescaleDB, servidor propio & N \\
Almacenamiento inmutable & Datos & Almacenamiento con redundancia de zona e inmutabilidad & N \\
Capa anticorrupción & Integración y eventos & Pasarela en el nodo de San Bernardo, junto al sistema contable & SB \\
Nodo de terminal & Servicios de negocio & Dos Karbon 430 en activo y en espera & GT \\
Búfer a bordo & Datos, versión reducida & eMMC del iWave G26I & DB \\
Réplica de recuperación & Todas & Brazil South, en espera & N2 \\
\end{tabla}
\notaTabla{Fuente: elaboración propia.}

La matriz muestra que las capas de servicios y de datos existen en tres lugares, nube, terminal y
camión, porque son las que deben seguir funcionando sin enlace.


% ----------------------------------------------------------------------------
\subsection{Catálogo de servicios contratados en la nube}\label{sec:4-catalogo}
% ----------------------------------------------------------------------------

Todos los servicios del catálogo están disponibles en Chile Central y en Brazil South. Se verificó en
la interfaz pública de precios de Azure, que solo lista un servicio en una región cuando está
desplegado ahí \parencite{ms_precios_api}. \transv{RT-03.05}{8} y el \art{16.3}{11} piden privilegiar
servicios administrados cuando reducen el riesgo operacional, y el catálogo los usa en todas las capas.
En este caso la preferencia tiene una razón concreta. El área de tecnología del mandante son nueve
personas para cinco terminales, dos talleres y una flota que no está en ningún lugar fijo
\caso*{capítulo 8}{20}. Con servicios administrados, los parches, la alta disponibilidad y los
respaldos de la plataforma los opera el proveedor de nube, y en San Bernardo y los terminales quedan
solo los equipos que el \art{16.4}{12} exige para operar sin enlace. La \tabref{4-servicios} declara
el nivel, la configuración y la subred de cada uno.

\begin{tabla}{Servicios de Azure en la región primaria, con nivel y configuración}{4-servicios}{L{36mm} X L{24mm}}
Servicio & Nivel y configuración & Subred \\ \midrule
Front Door & Premium, con cortafuegos de aplicación y protección volumétrica & Global, borde \\
API Management & Premium, con unidades en las tres zonas & Borde \\
IoT Hub & S1 con 2 unidades: 800.000 mensajes de 4 KB al día & Integración \\
Device Provisioning Service & Enrolamiento con certificado X.509 del elemento seguro & Integración \\
Event Hubs & Premium, protocolo Kafka, réplica geográfica asíncrona a Brazil South & Integración \\
Kubernetes Service (AKS) & Nodos en las tres zonas, con escalado automático & Aplicación \\
PostgreSQL Flexible transaccional & Alta disponibilidad entre zonas y réplica de lectura en Brazil South & Datos \\
PostgreSQL Flexible de series & Servidor propio con TimescaleDB, réplica en Brazil South & Datos \\
Almacenamiento de objetos & Redundancia de zona, inmutabilidad para evidencia y documentos & Datos \\
Key Vault & Premium, claves en módulo criptográfico & Datos \\
Azure Firewall & Premium, en la red central & Red central \\
Puerta de ExpressRoute & ErGw1AZ, redundante entre zonas & Red central \\
Puerta de VPN & VpnGw1AZ, redundante entre zonas & Red central \\
Bastion & Acceso administrativo sin dirección pública & Red central \\
Monitor y Log Analytics & Registros de todos los ambientes, 12 meses en línea y 24 en archivo & Gestión \\
\end{tabla}
\notaTabla{Fuente: elaboración propia sobre \textcite{ms_iothub_cuotas}, \textcite{ms_eh_georep} y \textcite{ms_pg_overview}.}

IoT Hub S1 con dos unidades admite 800.000 mensajes de 4~KB al día \parencite{ms_iothub_cuotas}, y la
sección~4.2.8 muestra que la operación normal con 430 camiones usa unos 330.000. La réplica asíncrona de
Event Hubs se configura con un retraso máximo, y el servicio frena la entrada si ese retraso se supera
\parencite{ms_eh_georep}, lo que acota la pérdida de telemetría ante un desastre. La retención de
Monitor sigue el \art{21.3}{15}, que exige doce meses en línea y veinticuatro en archivo para los
eventos de seguridad. Las series de tiempo van en un servidor separado del transaccional, de modo que
el peak de escritura de la reconexión masiva no atrase la réplica de los datos del despacho.


La infraestructura se define como código, versionada en el repositorio del mandante, porque
\transv{RT-03.03}{8} no admite infraestructura creada manualmente por consola. Todo recurso lleva
etiquetas de ambiente, centro de costo y responsable, con presupuestos y alertas de desviación, como
pide el \art{16.3}{11}. \transv{RT-03.07}{8} exige declarar la estrategia de reversibilidad. Los
contenedores y el motor relacional son portables. Los servicios de ingesta, mensajería y analítica no lo
son sin reescritura, aunque Event Hubs expone el protocolo Kafka y eso reduce el esfuerzo de migrar el
bus.

% ----------------------------------------------------------------------------
\subsection{Topología de red y segmentación}\label{sec:4-red}
% ----------------------------------------------------------------------------

La red en Azure es de tipo central y radial. Una red central por región concentra el cortafuegos de
Azure y las puertas de ExpressRoute y de VPN. Cada ambiente tiene su propia red, conectada a la
central, con cinco subredes: borde, integración, aplicación, datos y gestión. Solo la subred de borde se
expone a Internet, y las de aplicación y datos son privadas, como exige el \art{16.3}{11}. Los servicios
administrados se alcanzan por puntos de conexión privados dentro de su subred. La \tabref{4-cidr}
declara el plan de direcciones.

\begin{tabla}{Plan de direcciones de la red}{4-cidr}{L{46mm} L{30mm} X}
Red & Rango & Contenido \\ \midrule
Central de Chile Central & 10.10.0.0/22 & Cortafuegos, puertas de ExpressRoute y VPN, Bastion \\
Producción & 10.11.0.0/20 & Borde, integración, aplicación, datos y gestión, una /24 cada una, salvo aplicación con /22 para los nodos de AKS \\
Preproducción & 10.12.0.0/20 & Misma división que producción \\
QA & 10.13.0.0/20 & Misma división, con menos nodos \\
Desarrollo & 10.14.0.0/20 & Misma división, con menos nodos \\
Central de Brazil South & 10.20.0.0/22 & Cortafuegos y puerta de VPN \\
Recuperación en Brazil South & 10.21.0.0/20 & Misma división que producción \\
San Bernardo & 172.16.0.0/22 & Servidores, gestión, videovigilancia y acceso, en redes virtuales locales separadas \\
Terminales regionales & 172.16.4.0/22 & Una /24 por terminal \\
\end{tabla}
\notaTabla{Fuente: elaboración propia. Rangos privados que no se superponen entre sí.}

Los rangos no se superponen, de modo que cualquier sitio puede enrutar a cualquier red de Azure sin
traducción de direcciones. La \figref{4-azurefig} muestra la región primaria con sus redes, subredes y
servicios.

\figuraHorizontal{figuras/04-arquitectura/region-primaria.png}{Región primaria Azure Chile Central: red central, red de producción, subredes, servicios con su nivel y conexiones con San Bernardo y los terminales}{4-azurefig}

Todo el tráfico entre redes pasa por el cortafuegos de la red central, y cada subred tiene su grupo de
seguridad con reglas de denegación por defecto. Los camiones entran por IoT Hub con certificado mutuo,
los portales por Front Door y la puerta de enlace, y San Bernardo y los terminales por la red central.
Ningún componente de datos es alcanzable desde Internet. La administración entra por Bastion, sin
dirección pública en las máquinas.


% ----------------------------------------------------------------------------
\subsection{Ambientes, alta disponibilidad y respaldos}\label{sec:4-ambientes}
% ----------------------------------------------------------------------------

La solución tiene cinco ambientes y cada uno vive en su propia suscripción, como muestra la
\figref{4-ambientesfig}.

\figura{figuras/04-arquitectura/ambientes.png}{Ambientes de la solución, cada uno en su suscripción, y promoción desde desarrollo hasta producción}{4-ambientesfig}

Desarrollo y QA están en Chile Central con niveles de servicio menores y menos nodos. Preproducción
repite la arquitectura de producción a escala reducida, y ahí se ejecutan las pruebas de carga que
\transv{RT-09.06}{21} fija en 1,5 veces el peak declarado. Producción corre en las tres zonas de Chile
Central. El ambiente de recuperación está en Brazil South, en espera, y se describe en la
sección~4.3.2. Los cinco se levantan desde el mismo código de infraestructura, de modo que el de
recuperación es idéntico a producción por construcción.

\transv{RT-03.02}{8} exige al menos dos zonas para todo componente con requisito de alta
disponibilidad. En producción, AKS reparte sus nodos en tres zonas, PostgreSQL mantiene una copia en
espera en otra zona, el almacenamiento es de redundancia de zona y API Management tiene unidades en las
tres. En San Bernardo, los servidores, los firewalls y los switches van en pares, y en cada terminal los
dos Karbon 430 funcionan en activo y en espera.

Los respaldos siguen el esquema 3-2-1-1-0 de \transv{RT-07.09}{18}. Las tres copias son la base en
producción, su respaldo diario en Chile Central y la réplica en Brazil South. Los dos medios son el
motor de base de datos y el almacenamiento de objetos. La copia fuera de sitio es la de Brazil South.
La copia inmutable es el respaldo en almacenamiento con bloqueo de escritura por el período de
retención. Cero errores se verifica con una restauración de prueba mensual y con las dos conmutaciones
reales al año que pide \transv{RT-07.07}{17}. La custodia en medio físico de \transv{RT-06.26}{16} se
cubre con esa copia inmutable fuera de sitio, que el mismo requisito admite como solución más segura y
eficiente, porque la sala de San Bernardo es de sitio y no aloja el núcleo.

% ----------------------------------------------------------------------------
\subsection{Enlaces, puntos únicos de falla y conmutación}\label{sec:4-enlaces}
% ----------------------------------------------------------------------------

\transv{RT-03.17}{9} exige que el enlace entre el sitio on-premise y la nube sea redundante, con
caminos físicos y proveedores distintos, y con tiempo de conmutación declarado. La \tabref{4-enlaces}
declara cada enlace.

\begin{tabla}{Enlaces de la solución}{4-enlaces}{L{28mm} L{30mm} L{26mm} X L{25mm}}
Enlace & Medio principal & Capacidad & Respaldo & Conmutación \\ \midrule
San Bernardo y Azure & ExpressRoute en Santiago & 100 Mbit/s & VPN por el segundo proveedor, con otra ruta física & Hasta 1 minuto \\
Terminal y Azure & Enlace fijo con VPN & 10 Mbit/s o más & Router 5G con dos SIM & Bajo 1 minuto \\
Camión y Azure & LTE & 13 a 16 MB al mes & Iridium, mensajes de hasta 340 bytes & Inmediata para eventos críticos \\
Terceros y Azure & Plataformas de dos proveedores & Según cada proveedor & Ninguno & No aplica \\
\end{tabla}
\notaTabla{Fuente: elaboración propia sobre \textcite{ms_er_locations}, \textcite{ms_er_bfd} y \textcite{groundcontrol_edge}.}

ExpressRoute tiene punto de conexión en Santiago, en EdgeConnex SCL, con Cirion Technologies, Equinix y
PitChile como proveedores \parencite{ms_er_locations}. Santiago no ofrece la modalidad Metro ni la Local
de ExpressRoute, así que la redundancia la da la VPN por el segundo proveedor que San Bernardo ya tiene
\caso*{capítulo 5}{13}. Con detección bidireccional de fallas en el enlace privado, la caída se detecta
en menos de un segundo y la convergencia hacia la ruta de respaldo toma hasta un minuto
\parencite{ms_er_bfd}. La tabla de regiones de Microsoft asocia Santiago solo a Brazil South dentro de
Sudamérica, por lo que el proveedor debe confirmar que el circuito alcanza Chile Central con el nivel
Estándar. Los terminales conmutan en el propio RUTX50, que vigila el enlace fijo y pasa a la red 5G
\parencite{teltonika_rutx50}. Se configura con sondeos cada 10 segundos, de modo que la conmutación
ocurre en menos de un minuto. El enlace de los terceros no tiene respaldo porque depende de plataformas
que audIT no controla.

La cobertura móvil real de las rutas se caracteriza con mediciones en terreno, porque
\caso{RT-03.24}{31} declara que la disponibilidad informada por los operadores no es un antecedente
aceptable para el diseño. Esa campaña es una actividad de la Etapa~1 con costo y plazo, y produce dos
entregables con un solo recorrido: el mapa de sombras georreferenciado y el catálogo de lugares seguros
de detención que el criterio 28 necesita y que hoy nadie posee. La \tabref{4-spof} recorre los puntos
únicos de falla de la infraestructura y cómo se cubre cada uno.

\begin{tabla}{Puntos únicos de falla de la infraestructura y su mitigación}{4-spof}{L{42mm} X}
Punto de falla & Mitigación \\ \midrule
Enlace de San Bernardo & VPN por otro proveedor y otra ruta física \\
Enlace de un terminal & Router 5G con dos SIM \\
Nodo de un terminal & Segundo Karbon 430 en espera \\
Servidor de San Bernardo & Segundo R360 en alta disponibilidad \\
Firewall o switch de San Bernardo & Segundo equipo del par, en otro circuito \\
UPS de San Bernardo & Segunda UPS en la alimentación B \\
Grupo electrógeno & 66 minutos de UPS y reabastecimiento desde el estanque propio \\
Clima de San Bernardo & Segunda unidad, N+1 \\
Zona de Azure & Servicios en las tres zonas de Chile Central \\
Región de Azure & Brazil South en espera (sección 4.3.2) \\
Cobertura celular & Búfer de 288 horas y mensajería Iridium \\
Equipo a bordo & Reposición del 10 \% cambiada en el terminal \\
\end{tabla}
\notaTabla{Fuente: elaboración propia.}

Ningún punto de la tabla deja la operación detenida por sí solo. Subsisten dos puntos únicos que se
declaran como exige RT-02.11: el grupo electrógeno, que es uno solo porque la sala es de sitio y la
redundancia de generación es deseable y no obligatoria \transv*{RT-06.12}{15}, y el equipo a bordo,
que solo se cambia cuando el camión pasa por un terminal, cada seis días en promedio
\caso*{capítulo 10}{23}. Mientras tanto, el búfer guarda lo registrado.

\transv{RT-08.04}{18} pide fuentes de poder redundantes conectadas a circuitos distintos. Los servidores
R360 tienen dos fuentes, una en cada alimentación, y el FortiGate 90G admite un segundo adaptador, que
se compra \parencite{fortinet_fg90g}. El FortiSwitch 124F-POE y los equipos de gabinete tienen una sola
fuente, y en ellos la redundancia la da el par: cada equipo del par va en un circuito distinto.

% ----------------------------------------------------------------------------
\subsection{Dimensionamiento y capacidad}\label{sec:4-dimensionamiento}
% ----------------------------------------------------------------------------

El \caso{numeral 14.2}{30} exige estimar la volumetría del sistema y advierte que todo valor entregado
sin su derivación se evalúa como dimensionamiento no realizado. Advierte también que el problema de
esta operación es qué ocurre cuando trescientos camiones recuperan cobertura al mismo tiempo. La
estimación usa los supuestos de muestreo de la sección~4.2.9: posición cada 30~segundos en marcha y
cada 5~minutos detenido, telemetría del motor cada 60~segundos y 30 horas de marcha en 72 horas. La
\tabref{4-buffer} muestra lo que acumula un camión en 72 horas sin cobertura.

El perfil de ensayo es explícito y sintético: 30 h en marcha con posición cada 30 s y 42 h detenido cada 300 s producen \(3.600+504=4.104\) posiciones de 64 bytes; 30 h de motor cada 60 s producen 1.800 muestras de 160 bytes. Se propone un fixture de 45 eventos de jornada/conducción/esperas de 600 bytes, cinco documentos de 40.000 bytes y ocho fotos de 307.500 bytes. Los tamaños incluyen la serialización del fixture y se contrastan con firmware y documentos homologados; no son mediciones ni cantidades universales del caso. No se presupone compresión.

La \tabref{4-buffer} separa los componentes. Los eventos/documentos suman \(45\times600+5\times40.000=227.000\) bytes; todo el lote sin fotos suma 777.656 bytes.

\begin{tabla}{Perfil sintético serializado de un camión tras 72 h}{4-buffer}{X R{25mm} R{35mm}}
Componente & Cantidad & Bytes \\ \midrule
Posiciones & 4.104 & 262.656 \\
Motor & 1.800 & 288.000 \\
Eventos y documentos & 45 + 5 & 227.000 \\
Total sin fotos & 5.954 & 777.656 \\
Fotos adjuntas & 8 & 2.460.000 \\
Total de datos y fotos & Un lote & 3.237.656 \\
\end{tabla}
\notaTabla{Fuente: hipótesis analíticas y fixture de audIT; frecuencias S-02/S-03 y variantes CP-PERF-02/CP-HW-03 de T-17.}

Trescientos camiones producen 1.786.200 registros y 971.296.800 bytes incluidos adjuntos. Dividir ese volumen por veinte minutos requiere 6,48 Mbit/s útiles agregados, antes de overhead. Como presupuesto adicional se incorpora 25\,\% de transporte/metadatos y se reserva 20\,\% de la ventana para reconexión, reintentos, persistencia y conciliación: \(971.296.800\times1,25\times8/960=10,12\) Mbit/s. El porcentaje es una hipótesis declarada que se contrasta en pruebas; la garantía contractual solo se comprueba midiendo la confirmación íntegra de cada camión en hasta veinte minutos.

La \tabref{4-reconexion} distingue tasas de operaciones y cuotas. Los bloques de 4 KiB son una elección de empaquetado del ensayo y unidad de cuota diaria, no el tamaño máximo de un mensaje ni la tasa de operaciones \parencite{ms_iothub_cuotas}.

\begin{tabla}{Reconexión de 300 camiones: estimación y límites}{4-reconexion}{X R{31mm} L{48mm}}
Cálculo & Resultado & Alcance \\ \midrule
Paquetes de 4 KiB sin fotos & 57.000 & 300 por $\lceil777.656/4.096\rceil$ \\
Con presupuesto de 25\,\% & 71.400 & 238 paquetes por camión \\
Envíos a 100 operaciones/s & 11,9 min & Límite de dos unidades S1 \\
Nuevas conexiones a 100/s & Al menos 3 s & Solo admisión, sin reintentos \\
Inicio de 300 archivos agrupados & Al menos 90 s & 100/min por unidad; dos unidades \\
\end{tabla}
\notaTabla{Fuente: elaboración propia y documentación de Microsoft IoT Hub consultada el 7 de octubre de 2026.}

Estas estimaciones no prueban el tiempo final: admitir un archivo no lo transfiere, y aceptar un mensaje no confirma escritura ni conciliación. La transferencia de fotos comparte presupuesto de enlace; sus huellas deben compararse antes de cerrar el cronómetro. Colas, escritura, reintentos y carga normal se incluyen en CP-PERF-02. El ensayo aplica a cada camión y falla si uno excede veinte minutos o pierde registros; 288 h se ensaya separadamente como ampliación de ingeniería. La secuencia de ingesta, desacoplamiento y confirmación duradera se ilustra en la \figref{4-reconexionfig}.

\figura{figuras/04-arquitectura/reconexion.png}{Reconexión masiva: del búfer del camión a la base de datos, con los límites de cada tramo}{4-reconexionfig}

La figura sigue los cinco tramos. Cada equipo espera un tiempo aleatorio antes de reconectarse y
reintenta con retroceso exponencial, de modo que los 300 no golpean la puerta en el mismo segundo. Del
IoT Hub los mensajes pasan a Event Hubs, que desacopla la llegada del procesamiento: los consumidores en
AKS escriben en lotes a la base de series a su propio ritmo, y el mensaje queda retenido en el bus
aunque la base se atrase. La admisión de ingesta es el cuello de botella previsto por el cálculo de cuotas;
el ensayo debe contrastarlo con transferencia de fotos, procesamiento y escritura.
La figura identifica profundidad de cola y conciliación como indicadores, sin declarar una saturación medida.

En operación normal cada camión manda 768 mensajes diarios, uno por minuto en marcha y uno cada cinco
minutos detenido. Con los 430 camiones proyectados son unos 330.000 mensajes diarios, más 71.400 paquetes presupuestados de una
reconexión, dentro de los 800.000 de las dos unidades. El resto de lo que pide el numeral 14.2 está en
la \tabref{4-estimaciones}.

\begin{tabla}{Otras estimaciones del numeral 14.2}{4-estimaciones}{X R{26mm} L{52mm}}
Estimación & Valor & Cómo sale \\ \midrule
Eventos por segundo, promedio & 8,5 & 374 camiones con 1.968 registros diarios cada uno \\
Eventos por segundo en la reconexión & Cerca de 3.000 & 300 camiones con 5.954 registros en 10 minutos \\
Transacciones por segundo en la asignación & Cerca de 1 & Cota con los 374 camiones asignados en una hora y 10 consultas por asignación \\
Datos móviles por camión & 13 a 16 MB al mes & Muestreo de 30 segundos con sobrecarga de protocolo \\
Datos móviles de los equipos audIT & 2,4 a 2,9 GB al mes & 182 equipos por el valor anterior \\
Ancho de banda por terminal & 10 Mbit/s & 12,5 visitas diarias con imagen de 1 GB \\
Ancho de banda de San Bernardo & 52 Mbit/s & 80 usuarios internos a 0,5 Mbit/s, más continuidad, por el crecimiento a 390 usuarios \\
\end{tabla}
\notaTabla{Fuente: elaboración propia basada en la memoria de cálculo de ancho de banda y balance de cargas de San Bernardo.}

La carga fuerte de esta operación es la ingesta en la reconexión, unas 350 veces el promedio. La
asignación, que es la transacción crítica de \caso{RT-09.01}{32}, pesa poco en volumen: los 96.000
viajes al año son 263 al día. El ancho de banda de los terminales lo fija la actualización del equipo a
bordo. Con 62 visitas diarias a terminal en toda la red, calculadas con la flota completa como cota,
12,5 por terminal, y una imagen de 1~GB, se necesitan 3,5~Mbit/s repartidos en 8 horas y unos 9~Mbit/s
si cuatro camiones se actualizan en la misma hora de relevo. El ancho de banda de San Bernardo lo fijan
los usuarios de la casa matriz y de la torre, y por eso el circuito ExpressRoute es de 100~Mbit/s, el
siguiente escalón comercial sobre 52~Mbit/s.

La observación 75 pidió además el volumen transaccional anual y el de la evidencia que se almacena. La
\tabref{4-transaccional} los reúne.

\begin{tabla}{Volumen transaccional y de evidencia al año}{4-transaccional}{X R{30mm} R{30mm}}
Transacción & Hoy & A tres años \\ \midrule
Viajes & 96.000 & 118.000 \\
Documentos electrónicos de transporte & 128.000 & 157.000 \\
Abastecimientos de combustible & 74.000 & 90.000 \\
Pasadas de peaje & 620.000 & 760.000 \\
Órdenes de trabajo de taller & 5.200 & 6.100 \\
Liquidaciones a transportistas & 1.776 & 2.100 \\
Evidencia en almacenamiento inmutable & Cerca de 59 GB & Cerca de 67 GB \\
\end{tabla}
\notaTabla{Fuente: Caso, numeral 14.1, p. 29. Liquidaciones: 148 y 175 al mes por 12. Evidencia: documentos de transporte a 40 KB (5,1 GB hoy) más 2,46 MB de fotos cada 72 horas por equipo audIT (54 GB con 182 equipos y 61 GB con 204).}

La evidencia se retiene cinco años para la jornada y seis para los documentos de transporte
\caso*{RT-05.10}{31}, de modo que en los 56 meses del contrato se acumulan cerca de 290~GB, con lo que el
almacenamiento inmutable no condiciona el diseño. Los umbrales del numeral 9.1 transversal son exigibles
en producción medidos en el percentil 95 sobre la experiencia real de la persona usuaria, y en cada
fila prevalece el umbral más exigente entre ese numeral y el Caso. \transv{RT-09.03}{21} obliga a
soportar sin rediseño un crecimiento de tres veces la volumetría inicial, y la proyección del Caso llega
a 1,2 veces, de modo que esta oferta dimensiona sobre el requisito transversal. \transv{RT-09.05}{21}
pide identificar el primer cuello de botella. El modelo propone la admisión de ingesta durante
la reconexión masiva como limitación inicial; CP-PERF-02 debe verificarlo frente al enlace,
las fotos y la persistencia, midiendo profundidad de cola y retraso de procesamiento.
La ampliación de unidades de IoT Hub y de procesamiento se evalúa según el tramo medido, sin presuponer que elimine otro cuello de botella. La disponibilidad comprometida es de 99,9~\% mensual para los servicios
críticos, medida sobre la transacción de negocio de extremo a extremo \art*{20}{14}.


% ----------------------------------------------------------------------------
\subsection{Supuestos declarados y brechas}\label{sec:4-supuestos}
% ----------------------------------------------------------------------------

El \caso{numeral 14.2}{30} advierte que todo valor entregado sin su derivación se evalúa como
dimensionamiento no realizado. Los supuestos que sostienen las cifras de esta sección se declaran en la
\tabref{4-supuestostab} con la forma en que se cierran.

\begin{tabla}{Supuestos de audIT con su valor y su cierre}{4-supuestostab}{L{11mm} X L{32mm} L{36mm}}
ID & Supuesto & Valor & Cómo se cierra \\ \midrule
S-01 & Velocidad comercial media & 55 km/h & Medición en la Etapa 1 \\
S-02 & Horas de marcha dentro de 72 horas & 30 horas & Régimen de descansos del artículo 25 bis \\
S-03 & Tamaño del registro de posición & 64 bytes & Ajuste con el firmware del G26I \\
S-04 & Volumen acumulado a bordo tras 72 horas & 0,78 MB sin fotos y 2,46 MB de fotos & Derivado en la sección 4.2.8 \\
S-05 & Capacidad de almacenamiento a bordo & 2,52 GB asignados sobre 8 GB nominales & Derivada en la sección 4.2.2 \\
S-06 & Consumo mensual de datos por camión & 13 a 16 MB & Campaña de medición \\
S-07 & Imagen de actualización del equipo a bordo & 1 GB & Medición con el primer firmware \\
S-08 & Población con módem satelital & Los 182 equipos audIT & Decisión de diseño, no depende de la medición \\
S-09 & Parque telemático de los 192 camiones de terceros & Marca y modelo desconocidos & Clasificación contra el estándar de la sección 4.2.1 \\
S-10 & Consumo de cada servidor de San Bernardo & 300 W & Herramienta de planificación del fabricante \\
S-11 & Coeficiente del clima de precisión & 3,27 & Ficha del equipo a 50 Hz al comprar \\
S-12 & Tipología del recinto de San Bernardo & Sala técnica de sitio & Consulta al mandante \\
S-13 & Concurrencia simultánea en hora punta & 380 sesiones & Prueba de carga de RT-09.06 \\
S-14 & Registros de posición al año & Cerca de 120 millones & Derivada de S-01 y del muestreo \\
S-15 & Volumen almacenado de series & 30 a 40 GB al año & Modelo de datos del Subdocumento 5 \\
S-16 & Tráfico por usuario interno & 0,5 Mbit/s & Medición en la Etapa 1 \\
S-17 & Ventana de disponibilidad de las plataformas de terceros y de la autoridad tributaria & No declarada por las bases & Consulta al mandante \\
S-18 & Suscripción de Azure & A nombre del mandante & Reversibilidad de RT-03.07 \\
\end{tabla}
\notaTabla{Fuente: elaboración propia. S-01 a S-17 vienen del Informe 1 y se actualizaron con las decisiones de esta instancia.}

Cada supuesto tiene un mecanismo de cierre en la Etapa~1 o en la compra. Ninguno fija un valor que las
bases entreguen de otra forma.

Esta oferta reconoce cinco brechas que la ingeniería no cierra por sí sola.

\begin{enumerate}
\item La cobertura móvil no se supone. El Caso declara que la disponibilidad informada por los
operadores no es un antecedente aceptable para el diseño, y cualquier cifra sobre la calidad de la
señal en cada tramo antes de la medición en terreno sería una invención.
\item El parque telemático de los 192 camiones de terceros tampoco se supone. El Caso no entrega marca,
modelo, protocolo ni capacidad de almacenamiento de esos equipos. El entregable correcto es el
estándar mínimo de homologación de la sección~4.2.1, y el levantamiento es una actividad de la Etapa~1
con costo, plazo y dependencia de terceros.
\item El conductor que manejó otro camión sin dispositivo no deja rastro instrumental. Esa brecha se
cierra con responsabilidad contractual.
\item Las consultas del Artículo~43 están enviadas y no respondidas. Las interpretaciones que esta
oferta adopta sobre la sustitución de los módulos operativos de 2013, la autorización de los
fabricantes para leer la telemetría, el mecanismo digital de la red de estaciones de servicio y el
alcance del costo consolidado en 24 horas son propuestas de audIT, no respuestas del mandante.
\item El producto satelital evaluado en el Informe~1 guarda 40 horas de registro \parencite{webfleet_sat_cl},
frente a las 72 exigidas. Por eso esta oferta integra el mismo módem al G26I, que guarda 288 horas.
\end{enumerate}

% ============================================================================
\section{Data center}\label{sec:4-datacenter}
% ============================================================================

La solución usa dos regiones de Azure y dos tipos de sitio propio del mandante. El \transv{numeral
6.1}{14} obliga a declarar la tipología de cada sitio, advierte que sobredimensionar se castiga igual
que quedarse corto y fija para el recinto una disponibilidad de infraestructura de 99,95~\%. La
\tabref{4-tipologia} declara la tipología de cada sitio.

\begin{tabla}{Tipología declarada por sitio}{4-tipologia}{L{46mm} L{44mm} X}
Sitio & Tipología & Fundamento \\ \midrule
San Bernardo, 26 m² & Sala técnica secundaria o de sitio & Continuidad de la torre, no aloja el núcleo \\
Cuatro terminales regionales & Gabinete o borde operacional & Caso, RT-06.01, p. 32 \\
182 camiones con equipo audIT & On-premise distribuido & Caso, RT-06.01, p. 32 \\
\end{tabla}
\notaTabla{Fuente: FEP02, numeral 6.1, p. 14, y Caso, RT-06.01, p. 32.}

San Bernardo no aloja el núcleo, que está en la nube. Aloja lo que mantiene viva la operación cuando
cae el enlace, y por eso es una sala de sitio. A esa tipología el numeral 6.1 le aplica los requisitos
de energía, climatización, control de acceso, detección de incendio y monitoreo, dimensionados al sitio.

Las bases fijan dos plazos distintos para operar sin enlace. \caso{RT-03.10}{31} pide que los
terminales operen 12 horas sin enlace hacia el exterior, y el \art{16.4}{12} exige que el componente
on-premise opere de forma autónoma degradada al menos 24 horas continuas, «en ningún caso inferior». Esta
oferta diseña San Bernardo y los terminales para 24 horas, que cumple ambos textos, y deja la diferencia
declarada para que el mandante la aclare.

% ----------------------------------------------------------------------------
\subsection{Especificaciones Data Center Primaria}\label{sec:4-dc-primario}
% ----------------------------------------------------------------------------

La región primaria es Azure Chile Central, en Santiago, con tres zonas de disponibilidad
\parencite{ms_regions}. El proveedor cumple \transv{RT-03.01}{8}, que exige región en Chile o en
Sudamérica, y los componentes con alta disponibilidad se reparten en las tres zonas, como pide
\transv{RT-03.02}{8}. La residencia de los datos queda en el país y sujeta a la aprobación del mandante,
como dispone el \art{23}{17}. Chile Central no tiene región pareja \parencite{ms_regions}. Eso importa
para la recuperación: los servicios que Azure replica solos hacia la región pareja no tienen adónde
replicar, y la copia a Brazil South se diseña servicio por servicio en la sección~4.3.2.

La infraestructura on-premise asociada está en la sala de San Bernardo, que es casa matriz, sede de la
torre 24x7 y terminal principal \caso*{capítulo 3}{8}. Hoy la sala mide 26~m², se habilitó en 2013,
tiene climatización split, 20 minutos de alimentación ininterrumpida y acceso por credencial, y el Caso
declara que no cumple el capítulo 6 de las Transversales \caso*{capítulo 5}{13}. La
\tabref{4-habilitacion} compara cada exigencia con lo que hay y con lo que se instala.

\begin{tabla}{Habilitación de la sala de San Bernardo}{4-habilitacion}{L{20mm} L{38mm} L{19mm} X L{19mm}}
Ámbito & Exigencia & Hoy & Solución & Fuente \\ \midrule
Energía & UPS de 30 minutos a plena carga & 20 minutos & Dos APC SRT3000 en alimentación A y B con batería externa & RT-06.07 \\
Energía & Generación de 24 horas con estanque y contrato & No existe & Himoinsa HYW-8 T5 S5 con estanque de 100 L & RT-06.08 \\
Energía & Instalación propia y puesta a tierra según NCh Elec. 2777 & Por verificar & Tablero propio con transferencia automática & RT-06.09 \\
Energía & Carga, factor de potencia y PUE declarados & No declarados & 1,94 kW, factor 0,95 y PUE de 1,95 & RT-06.11 \\
Clima & Precisión en N+1 & Split & Dos Liebert Mini-Mate2 de 3,55 kW & RT-06.13 \\
Clima & Temperatura, humedad y agua en línea & No existe & Controlador del clima y Liebert LT460 & RT-06.14 \\
Incendio & Detección por aspiración láser & No existe & VESDA VLF-250 & RT-06.16 \\
Incendio & Extinción con agente limpio aprobado por UL & No existe & Kidde Fluoro-K, FK-5-1-12 & RT-06.17 \\
Acceso & Biometría facial con AFIS de respaldo & Credencial & Dos Suprema FaceStation F2 en esclusa & RT-06.20 \\
Acceso & Videovigilancia de 30 días & No existe & Cuatro cámaras Axis y grabador de 4 TB & RT-06.24 \\
Redes & Rutas físicas distintas & Dos proveedores & ExpressRoute y VPN por ductos separados & RT-06.32 \\
\end{tabla}
\notaTabla{Fuente: FEP02, capítulo 6, pp. 14 a 17, y Caso, capítulo 5, p. 13.}

La brecha está en energía, clima, incendio y acceso, que son los ámbitos que el numeral 6.1 exige para
una sala de sitio. El agente de extinción es el FK-5-1-12. El Informe~1 proponía el Novec~1230, que es
ese mismo compuesto, pero 3M dejó de fabricarlo a fines de 2025 \parencite{fike_novec}. Kidde lo
ofrece como Fluoro-K, con aprobación UL y FM \parencite{kidde_fluorok}, lo que cumple
\transv{RT-06.17}{15}. La detección y la extinción se integran al monitoreo, que avisa al centro de
operación y a la contraparte del mandante \transv*{RT-06.19}{16}, y la sala suma dos extintores
portátiles de dióxido de carbono \transv*{RT-06.18}{15}.

\paragraph{Balance eléctrico y térmico} \transv{RT-06.11}{15} pide declarar la carga proyectada en kW,
el factor de potencia y el PUE. La \tabref{4-balance} suma la carga con los consumos de ficha de cada
equipo y los supuestos declarados.

\begin{tabla}{Balance eléctrico de la sala de San Bernardo}{4-balance}{X R{16mm} L{58mm}}
Concepto & W & Base \\ \midrule
Dos servidores Dell R360 & 600 & 300 W cada uno, supuesto S-10 \\
Dos firewalls FortiGate 90G & 40 & 19,9 W cada uno, ficha del fabricante \\
Dos switches FortiSwitch 124F-POE & 105 & 237,4 W máximo menos 185 W de PoE, ficha \\
Cuatro cámaras de videovigilancia & 52 & 12,95 W de la clase 3 de PoE cada una \\
Acceso, grabador, detección y sensores & 150 & Supuesto de diseño \\
Equipos de los dos proveedores de enlace & 50 & 25 W cada uno, supuesto de diseño \\
Carga de TI & 996 & Suma de las filas anteriores \\
Pérdidas de las dos UPS & 100 & 7,6 \% a plena carga según la ficha \\
Iluminación & 208 & 8 W/m² en 26 m² \\
Clima, una unidad en servicio & 640 & 1,6 kW de calor con COP 3,27, más 149 W de ventilador \\
Carga total & 1.944 & PUE de 1,95 \\
\end{tabla}
\notaTabla{Fuente: \textcite{fortinet_fg90g}, \textcite{fortinet_fs124f}, \textcite{apc_srt3000}, \textcite{stulz_minispace}, \textcite{vertiv_minimate}. Memoria de cálculo térmico y eléctrico según estándares ASHRAE TC 9.9 y TIA-942.}

La carga de TI es de cerca de 1~kW y la total de 1,94~kW, con un PUE estimado de 1,95. El PUE es alto
porque en una sala de 26~m² el clima y la iluminación pesan mucho frente a una carga de TI chica. El
calor a retirar, 1,6~kW, cabe en una sola unidad de 3,55~kW, y la segunda queda en reserva para cumplir
el N+1 de \transv{RT-06.13}{15}. El coeficiente de 3,27 sale de la ficha de una unidad de precisión de
Stulz, 6,6~kW de capacidad con 2.020~W de compresor \parencite{stulz_minispace}, y se reemplaza por el del
equipo comprado. El factor de potencia de diseño es 0,95, porque los servidores y las UPS de doble
conversión corrigen el factor en su entrada, y se mide en la revisión semestral de
\transv{RT-06.10}{15}.

La UPS se dimensiona con la carga de TI. Cada APC SRT3000 tiene 505~VAh internos, y con una batería
externa llega a 1.515~VAh \parencite{apc_srt3000}. Con un rendimiento del inversor de 0,9 y la batería al
80~\% de su capacidad al final de su vida, una sola UPS sostiene toda la carga de TI durante 66 minutos
(1.515 × 0,9 × 0,8 / 996 W), el doble de los 30 que pide \transv{RT-06.07}{15}. Sin la batería externa
serían 22 minutos, y por eso se compra.

El grupo electrógeno toma la carga total más la recarga de las dos UPS, 168~W cada una: 2,28~kW, o
2,85~kVA a factor 0,8. El Himoinsa HYW-8 T5 S5 entrega 8,8~kVA y queda al 32~\% \parencite{himoinsa_hyw8}.
Un grupo de 15~kVA quedaría bajo el 30~\%, carga con la que un motor diésel acumula combustible sin
quemar. A plena carga consume 2,68 litros por hora, 64,3 litros en 24 horas, y el estanque opcional de
100 litros cumple las 24 horas de \transv{RT-06.08}{15} con margen. El reabastecimiento se hace desde el
estanque de combustible propio que la compañía tiene en San Bernardo \caso*{numeral 2.2}{6}, con un
contrato de respaldo con el proveedor de combustible.

La videovigilancia guarda 30 días en línea, como pide \transv{RT-06.24}{16}. Cuatro cámaras Axis
M3215-LVE a 2~Mbit/s en H.265 generan 2,59~TB en 30 días, y el grabador Axis S3008 Mk II de 4~TB los
cubre \parencite{axis_m3215,axis_s3008}. Las grabaciones anteriores se copian al almacenamiento inmutable
de Azure, lo que suma 8~Mbit/s al enlace de San Bernardo, dentro de su capacidad de 100~Mbit/s.

\paragraph{Plano de la sala} \transv{RT-06.03}{14} pide el plano de distribución con las zonas de
generadores, baterías, climatización, servidores, comunicaciones, trabajo y respaldo. El numeral 6.1 lo
exige a la sala principal, y esta oferta lo entrega igual porque la revisión del Informe~1 lo pidió. La
\figref{4-planofig} lo muestra sobre una planta supuesta de 6,5 por 4~metros, que suma los 26~m²
del Caso.

\figura{figuras/04-arquitectura/sala-san-bernardo.png}{Plano de la sala de San Bernardo con sus zonas, a proporción de 26 m²}{4-planofig}

Se entra por una esclusa con dos lectores biométricos, uno en cada puerta, que no deja pasar a más de
una persona a la vez \transv*{RT-06.23}{16}. El rack de servidores y el de comunicaciones van separados,
como pide \transv{RT-06.05}{15}, y cada uno es de 24 unidades. El de servidores ocupa cerca de 12
unidades con los dos R360, las dos UPS y sus baterías externas, y el de comunicaciones cerca de 8 con
los firewalls, los switches, los paneles de conexión y los equipos de los proveedores, de modo que ambos
quedan con la mitad o más libre para crecer. La \figref{4-racksfig} muestra la elevación de los dos racks, con las unidades que ocupa cada equipo y las que quedan libres.

\figura{figuras/04-arquitectura/racks-san-bernardo.png}{Elevación frontal de los racks de servidores y comunicaciones en la sala técnica de San Bernardo}{4-racksfig}

La disponibilidad de más del 50\% de espacio vertical asegura la escalabilidad y adición modular de hardware durante todo el período de concesión. Las dos unidades de clima van en muros opuestos, el cilindro
de FK-5-1-12 junto a la esclusa, el puesto de trabajo con la consola en la zona de trabajo y los
repuestos de sala en la zona de respaldo. El grupo electrógeno y su tablero de transferencia quedan en
el patio, fuera de la sala, y los dos enlaces entran por ductos separados, como pide
\transv{RT-06.32}{17}. Las medidas de detalle se confirman en el levantamiento de la Etapa~1, porque el
Caso solo entrega la superficie.

\paragraph{Gabinetes de los terminales regionales} Cada terminal regional lleva el gabinete de la
\figref{4-terminalfig}.

\figura{figuras/04-arquitectura/gabinete-terminal.png}{Gabinete de un terminal regional, con sus enlaces, el punto de acceso de patio y el lector de portería}{4-terminalfig}

Dentro del gabinete van dos OnLogic Karbon 430 en activo y en espera, que funcionan sin ventilador de
$-40$ a $+70$~°C con alimentación de 9 a 48~V y consumen cerca de 35~W cada uno \parencite{onlogic_k430}.
Con dos equipos el terminal no queda con un punto único de falla. El router Teltonika RUTX50, 5G con dos
SIM y conmutación automática, opera de $-40$ a $+75$~°C y consume menos de 18~W
\parencite{teltonika_rutx50}. El switch Teltonika TSW202, de riel DIN y con PoE, conecta los equipos
\parencite{teltonika_tsw202}. Fuera del gabinete, el FortiAP 234G, IP67 y de $-50$ a $+70$~°C, cubre el
patio con Wi-Fi para actualizar los camiones \parencite{fortinet_fap}, y el Minew G1 lee en la portería
las balizas de los semirremolques \parencite{minew_g1}. El gabinete suma 152~W, y la UPS APC SRT1000XLI,
con 293~VAh \parencite{apc_srt1000}, lo sostiene 83 minutos al final de la vida de su batería. Sin enlace,
el terminal sigue registrando salidas y llegadas, verifica la aptitud contra la copia local de vigencias
y jornada, y actualiza los camiones con la imagen que guarda, durante las 24 horas que fija el
\art{16.4}{12}.

El monitoreo de San Bernardo y de los terminales se integra a la misma plataforma de observabilidad que
la nube, con alertas unificadas, como pide \transv{RT-03.16}{9}.

% ----------------------------------------------------------------------------
\subsection{Especificaciones Data Center Secundario}\label{sec:4-dc-secundario}
% ----------------------------------------------------------------------------

La región secundaria corresponde a Azure Brazil South (São Paulo), provista con tres zonas
de disponibilidad \parencite{ms_regions}. Santiago y São Paulo están a unos 2.586~km en línea recta, en
otro país, con otro sistema eléctrico y lejos de la zona de subducción que concentra los grandes
sismos de Chile, lo que responde a \transv{RT-07.02}{17}. San Bernardo, en la misma área metropolitana
que la región primaria, no puede ser el sitio de recuperación porque comparte con ella el riesgo
sísmico y la red eléctrica. Microsoft no publica la latencia entre las dos regiones
\parencite{ms_latencia}. El mínimo físico de ida y vuelta por fibra es de unos 26~ms (2.586~km × 2 /
200.000~km/s), y la latencia real se mide en la Etapa~1.

La modalidad es activo-pasivo, como permite \transv{RT-07.01}{17}. Una copia activa en Brazil South
duplicaría el costo de operación para mejorar un RTO que el activo-pasivo ya cumple, y agregaría la
complejidad de escribir en dos regiones a la vez. Brazil South queda en espera con réplicas asíncronas y
se promueve ante un desastre. Los objetivos son los de \transv{RT-07.04}{17}: 4 horas para recuperar el
servicio y no más de 15 minutos de datos perdidos. La \tabref{4-replicas} declara cómo se copia cada
servicio.

\begin{tabla}{Cómo llega cada servicio a Brazil South}{4-replicas}{L{34mm} X L{30mm} X}
Servicio & Mecanismo & Redundancia allí & Efecto en el RPO \\ \midrule
PostgreSQL transaccional & Réplica de lectura asíncrona que se promueve a primaria & Local o de zona & Retraso de segundos a minutos, con alerta bajo 15 minutos \\
PostgreSQL de series & Réplica de lectura asíncrona propia & Local o de zona & El peak de la reconexión no atrasa la transaccional \\
Event Hubs & Réplica geográfica asíncrona con retraso máximo de 10 minutos & Premium & Pérdida acotada a 10 minutos, que el búfer a bordo repone \\
Almacenamiento de objetos & Replicación de objetos a una cuenta de Brazil South & Local o de zona & Documentos copiados en minutos \\
AKS y servicios & Mismo código de infraestructura, clúster en espera & Tres zonas & Afecta el RTO, no el RPO \\
Key Vault & Bóveda propia con las claves respaldadas & Del servicio & Afecta el RTO, no el RPO \\
\end{tabla}
\notaTabla{Fuente: elaboración propia sobre \textcite{ms_pg_replicas}, \textcite{ms_eh_georep} y \textcite{ms_regions}.}

Ninguna copia en Brazil South usa redundancia geográfica. La pareja de Brazil South es South Central US,
en Texas \parencite{ms_regions}, y una redundancia geográfica llevaría los datos fuera de Sudamérica.
Microsoft advierte que el retraso de una réplica asíncrona de PostgreSQL suele ser de segundos a minutos
y puede llegar a horas con escritura intensa \parencite{ms_pg_replicas}. Por eso la telemetría tiene tres
defensas. Va a un servidor propio, de modo que su peak no atrasa la réplica transaccional. Pasa antes por
Event Hubs, cuya réplica frena la entrada si el retraso supera los 10 minutos \parencite{ms_eh_georep}. Y
el equipo a bordo conserva lo enviado durante 72 horas, así que tras una conmutación reenvía el tramo que
la réplica pudo no alcanzar, y la clave de idempotencia de siete días descarta los duplicados.

La \figref{4-drfig} muestra la recuperación entre regiones y el eje separado de continuidad.

\figuraHorizontal{figuras/04-arquitectura/recuperacion.png}{Recuperación ante desastres entre Chile Central y Brazil South, y eje separado de continuidad en San Bernardo, los terminales y los camiones}{4-drfig}

La conmutación sigue los cuatro pasos de la parte superior de la figura. El monitoreo detecta la caída y
alerta con un criterio declarado: la región primaria no responde durante 15 minutos, el plazo que
Microsoft sugiere como referencia para promover una región secundaria \parencite{ms_eh_georep}. La guardia
confirma que es un desastre y no una falla pasajera. Un procedimiento automatizado promueve las
réplicas y despliega los servicios en Brazil South. La operación pasa a la región secundaria. La
confirmación humana protege contra una conmutación innecesaria, como valora \transv{RT-07.08}{17}, y el
procedimiento queda documentado y automatizado como exige \transv{RT-07.05}{17}. La vuelta a Chile
Central sigue un procedimiento igual, con conciliación de lo generado durante la contingencia
\transv*{RT-07.06}{17}, y la conmutación se prueba de verdad dos veces al año
\transv*{RT-07.07}{17}.


La parte inferior de la figura es el otro eje. La recuperación y la continuidad son ejes separados:
Brazil South recupera la nube si cae la región primaria, y San Bernardo, los terminales y los camiones
mantienen la operación cuando lo que cae es el enlace con la nube. La \tabref{4-ejes} los separa.

\begin{tabla}{Los dos ejes de continuidad}{4-ejes}{L{40mm} L{40mm} X}
Eje & Primario y secundario & Fundamento \\ \midrule
Recuperación ante desastres & Azure Chile Central y Azure Brazil South & Distancia que evita compartir el evento (RT-07.02), con RTO de 4 horas y RPO de 15 minutos (RT-07.04) \\
Continuidad operacional & Nube, San Bernardo, terminales y camiones & RT-21.06 del Caso clasifica en severidad máxima todo incidente que impida asignar un viaje, emitir un documento de transporte o recibir un evento de emergencia \\
\end{tabla}
\notaTabla{Fuente: FEP02, RT-07.02 y RT-07.04, p. 17, y Caso, RT-21.06, p. 34.}

\transv{RT-03.13}{9} obliga a declarar qué funciones no estarán disponibles sin enlace y qué
procedimiento manual las suple, y evalúa como observación grave la ausencia de esta declaración. El
mismo código en el Caso regula otra materia, la sincronización en 20 minutos tras la reconexión, que la
sección~4.2.8 dimensiona. La \tabref{4-offline} declara las funciones.

\begin{tabla}{Disponibilidad de funciones sin enlace y procedimiento supletorio}{4-offline}{L{52mm} L{34mm} X}
Función & Sin enlace & Procedimiento supletorio \\ \midrule
Registro de posición, jornada y eventos & Disponible & No aplica \\
Evaluación de geocerca de llegada y salida & Disponible & No aplica \\
Alerta de jornada próxima a agotarse & Disponible & No aplica \\
Botón de emergencia & Disponible por satélite en los 182 equipos audIT & Protocolo telefónico con la torre para los camiones de terceros \\
Asignación de un viaje no precargado & No disponible & Autorización de la torre, registrada y reconciliada al recuperar enlace \\
Verificación bloqueante con dato fresco & Parcial, contra copia local & Regla de excepción con rol nominado y registro auditable \\
Consulta de liquidación por el transportista & No disponible & Portal, al recuperar enlace \\
Notificación en tiempo real al cliente & No disponible & Aviso diferido \\
Actualización del equipo a bordo & Solo en terminal & Con la imagen que guarda el gabinete \\
\end{tabla}
\notaTabla{Fuente: elaboración propia sobre FEP02, RT-03.13, p. 9.}

La emisión del documento electrónico de transporte sin cobertura no está en la tabla porque su mecanismo
lo define la sección~4.1, con el sistema contable como único emisor.

Los formularios de este subdocumento van en archivo propio: Formulario T-11 (\texttt{AUDIT-Formulario-T-11.pdf}).
